\chapter{采样式 MPC：路径积分理论与 MPPI}
\label{ch:mppi}

% ════════════════════════════════════════════════════════════════
%  采样式/随机 MPC 新范式章：全吸收。
%  融合 Tutorial 04_移动机器人规控/20_采样式MPC 的三篇笔记：
%   20_路径积分控制理论.md（Kappen 指数变换/Feynman-Kac/自由能-KL/温度/LMDP）
%   30_MPPI核心算法与GPU.md（权重公式工程细节，取理论相关部分）
%   50_CEM家族与统一视角.md（CEM/iCEM/Tsallis/Predictive Sampling/CoVO-MPC 统一框架）
%  导论 planning_intro 未覆盖随机/采样 MPC，本章为新范式；与其搜索/采样/轨迹优化基础以 \cref 桥接、不重复。
%  源 md 由 AI 撰写、多轮审仍可能有错/幻觉，存疑处就地 \pz/\rebuilt 并入 claimsToVerify。
%  教学层遵循 v5.0 理论教学规范。
% ════════════════════════════════════════════════════════════════

\begin{practice}[前置自测]
开始之前，先试答下列问题。若已能流畅作答，可只速读本章表格与速查卡；若卡住，正是本章要为你解决的。
\begin{enumerate}
  \item 写出有限时域\emph{随机}最优控制的 Hamilton--Jacobi--Bellman（HJB）方程。它对一般非线性系统为何没有解析解？方程里哪一项是非线性的？
  \item 什么是指数变换 / Cole--Hopf 变换？你在别处是否见过“取对数把乘法变加法、把 $\min$ 变成 $\log\sum\exp$”这种把非线性运算线性化的手法？
  \item 写出 KL 散度 $\mathrm{KL}(P\,\|\,Q)$ 的定义式。一个温度系数 $\lambda$ 调大、调小，分别让一个由代价决定的分布更“贪心”还是更“随机”？
  \item 什么是重要性采样？已知只能从分布 $Q$ 采样，如何估计另一分布 $P$ 下的期望 $\mathbb{E}_P[h(x)]$？
  \item 一个非线性 SDE 的最优控制，能不能不解任何偏微分方程、只靠“撒一批随机轨迹、做加权平均”得到？若能，权重该长什么样？
\end{enumerate}
\end{practice}

\section*{本章目标}
学完本章，你应当能够：
\begin{enumerate}
  \item \textbf{推导}指数变换 $\psi=e^{-J/\lambda}$ 如何在 Kappen 条件 $\boldsymbol{\Sigma}=\lambda R^{-1}$ 下把非线性 HJB 偏微分方程\emph{整体}线性化，并说清该条件的物理含义；
  \item \textbf{用} Feynman--Kac 公式把线性 backward PDE 写成“无控扩散过程下的路径积分期望”，并解释为何这意味着可用蒙特卡洛估计最优控制、无需从终点反扫；
  \item \textbf{从}信息论的自由能--KL 对偶重新推导路径积分控制，说清它如何去掉 Kappen 的控制仿射与二次控制代价限制，对任意黑箱代价与动力学成立；
  \item \textbf{独立推导} MPPI 权重公式 $w_k=\exp(-(S_k-\rho)/\lambda)/\sum_j\exp(-(S_j-\rho)/\lambda)$，解释 $\rho$ 减法的数值意义与高斯密度、控制代价为何在权重里消去；
  \item \textbf{论证} MPPI 与 REINFORCE 的数学同构——MPPI 是零阶策略梯度的闭式特例，并理解“无学习率”来自分布空间的自然梯度 / 镜像下降；
  \item \textbf{解释}温度 $\lambda$ 如何在“贪心取 $\arg\min$”与“一视同仁地平均”之间连续插值，及其同时是 softmax 锐度、噪声温度、风险态度的三重身份；
  \item \textbf{说出}路径积分在离散时间的镜像——Todorov 线性可解 MDP（$z=MPz$），并解释它独有的“最优策略可线性叠加”的可组合性；
  \item \textbf{把} MPPI 定位进采样式 MPC 家族，用统一的“采样—赋权—更新”骨架解释 CEM、iCEM、Tsallis、Predictive Sampling、CoVO-MPC 的异同；
  \item \textbf{用}接收水平控制把单次开环序列变成闭环反馈，并理解协方差 $\boldsymbol{\Sigma}$（撒哪里）与温度 $\lambda$（信哪条）是两个正交的旋钮。
\end{enumerate}

\subsection*{本章知识导航}
本章是一条“\emph{把非线性随机最优控制，化为撒一批轨迹做指数加权平均}”的主线。结构如下：

\begin{center}
\small
\begin{tikzpicture}[
  node distance=4mm and 7mm, >=Stealth,
  blk/.style={draw, rounded corners, align=center, font=\small, inner sep=3pt, fill=main!6, draw=main!60!black},
  grp/.style={draw, rounded corners, align=center, font=\small, inner sep=3pt, fill=second!6, draw=second!60!black},
  uni/.style={draw, rounded corners, align=center, font=\small, inner sep=3pt, fill=third!8, draw=third!60!black},
  ln/.style={->, thick, gray!60!black}]
\node[blk] (hjb) {非线性 HJB\\（随机最优控制）};
\node[blk, right=of hjb] (exp) {指数变换\\ $\psi=e^{-J/\lambda}$};
\node[blk, right=of exp] (fk) {Feynman--Kac\\ 路径积分期望};
\node[grp, below=8mm of hjb] (lmdp) {离散镜像\\ LMDP $z=MPz$};
\node[grp, right=of lmdp] (fe) {自由能--KL 对偶\\ 吉布斯分布};
\node[grp, right=of fe] (w) {MPPI 权重\\ $w_k\propto e^{-S_k/\lambda}$};
\node[uni, below=8mm of lmdp] (lam) {温度 $\lambda$\\ 软$\min$/风险};
\node[uni, right=of lam] (rein) {同构\\ REINFORCE/SAC};
\node[uni, right=of rein] (fam) {家族统一\\ CEM/iCEM/Tsallis/CoVO};
\draw[ln] (hjb)--(exp); \draw[ln] (exp)--(fk);
\draw[ln] (exp.south) -- ++(0,-4mm) -| (fe.north);
\draw[ln] (fk.south) -- ++(0,-4mm) -| (w.north);
\draw[ln] (fe)--(w);
\draw[ln] (hjb.south)--(lmdp.north);
\draw[ln] (w.south) -- ++(0,-4mm) -| (fam.north);
\draw[ln] (lam)--(rein);
\end{tikzpicture}
\end{center}

\noindent\textbf{推荐路径}：\cref{sec:mppi-hjb,sec:mppi-fk} 是\emph{理论地基}（指数变换 + Feynman--Kac，把“解 PDE”变成“撒点”），任何读者必读；\cref{sec:mppi-freeenergy} 是\emph{算法落地}的核心（自由能--KL 对偶给出权重公式，并把适用范围从 LQ 特例推广到任意黑箱），是全章篇幅最大、最承重的一节；\cref{sec:mppi-temperature}（温度 $\lambda$）解剖唯一的自由参数；\cref{sec:mppi-algorithm}（接收水平 MPPI 算法）把公式落成可部署形态；\cref{sec:mppi-family}（家族统一）把 MPPI 放回采样式 MPC 全景，与 CEM 谱系打通。带 LMDP、控制即推断、风险敏感标记的小节为进阶——其中 LMDP（\cref{sec:mppi-lmdp}）是离散侧的权威镜像、控制即推断（\cref{sec:mppi-inference}）把 MPPI/SAC/REINFORCE 收束到一根主轴。

\subsection*{前置知识桥接}
本章是 \cref{ch:planning}（运动规划导论）之后的\emph{新范式}：导论的搜索式（A$^*$、Dijkstra）、采样式（RRT/PRM）、优化式（CHOMP、minimum-snap，见 \cref{sec:plan-traj}）都假定一个\emph{确定性}的、可显式建模的世界，且多数把规划与底层控制分离。本章处理的是\emph{随机动力学下的在线最优控制}——把规划与反馈控制合一，每个控制周期重新求解一个随机最优控制问题。它与导论的接口有二：其一，本章的代价泛函 $S(\tau)=\phi(x_T)+\sum_t V(x_t)\Delta t$ 与 \cref{def:plan-trajopt} 轨迹优化的泛函同形，只是这里允许 $V$ 不可微、不连续（黑箱）；其二，本章反复出现的“接收水平 / 滚动时域”机制（\cref{sec:mppi-algorithm}）是所有 MPC 的共性，与导论的级联流水线“全局搜索 + 局部优化”互补——MPPI 常作为那个“局部优化”环节的随机版替代。本章另需 \cref{ch:lie} 不要求，但要求读者熟悉动态规划的 Bellman 递推（\cref{thm:plan-bellman}）、KL 散度与重要性采样、Itô 随机微积分的一句话直觉。

\subsection*{如果跳过本章会怎样}
\begin{itemize}
  \item 你会把 MPPI 权重 $\exp(-S/\lambda)$ 当成一个“玄学配方”——不知道 $\lambda$ 该取多大、为什么是指数而非别的函数、为什么减去 $\rho=\min_k S_k$，于是调参全靠试，换个机器人就从零撞运气；
  \item 你会看不出 MPPI 与你已会的 REINFORCE / SAC 是同一件事，于是 RL 那套调参直觉（reward shaping、温度退火、方差缩减）在 MPPI 里全部白白浪费；
  \item 你会把 MPPI、CEM、iCEM、Predictive Sampling 当成几门要分别死磕的孤立手艺，记不住、也看不出它们其实是同一个“采样—赋权—更新”框架在不同权重函数下的实例。
\end{itemize}

\begin{table}[H]\centering\small
\caption{本章阅读方式建议}
\label{tab:mppi-reading}
\begin{tabular}{lll}
\toprule
阅读方式 & 时间 & 适合谁 \\
\midrule
精读（含全部推导与练习） & 6--8 小时 & 首次系统学路径积分控制、要做 MPPI 实现/研究 \\
速读（跳过冗长推导细节） & 1.5--2 小时 & 已了解 MPPI、想补理论根与家族脉络 \\
速查（只看小结的符号表 + 定理速查表 + 算法清单） & 20 分钟 & 写代码、选采样式控制器时回来查 \\
\bottomrule
\end{tabular}
\end{table}

\begin{note}
\textbf{本章记号与约定.}\;
状态 $x\in\mathbb{R}^n$、控制 $u\in\mathbb{R}^m$；漂移 $f(x)$、控制矩阵 $G(x)$、噪声协方差 $\boldsymbol{\Sigma}$、控制代价权重 $R\succ0$、温度 $\lambda>0$。
最优代价 / 值函数 / cost-to-go 记 $J(x,t)$；desirability（合意度）$\psi=e^{-J/\lambda}$（连续时间）、$z=e^{-v/\lambda}$（离散时间）。
一条轨迹（时间序列）记 $\tau$、其总代价 $S(\tau)=\phi(x_T)+\sum_t V(x_t,t)\Delta t$（\textbf{标量}，允许黑箱/不可微）；运行代价 $V$、终端代价 $\phi$。
$K$ 条 rollout 的噪声序列 $\{\boldsymbol{\varepsilon}_k\}$、总代价 $\{S_k\}$、归一化权重 $\{w_k\}$、控制序列 $U=\{u_0,\dots,u_{H-1}\}$、时域 $H$（或 $T$，与连续时间终端时刻同名，由上下文区分）。
\textbf{字母隔离声明}：本章 $\lambda$ 指\emph{温度}（非特征值/波长）、$\boldsymbol{\Sigma}$ 指\emph{噪声/采样协方差}、$\rho$ 指数值稳定基准 $\min_k S_k$（\cref{sec:mppi-family} 中 CoVO-MPC 处另用 $\rho(\boldsymbol{\Sigma})$ 表收缩率，就地声明）、$r$ 指 Tsallis 变形指数参数；均仅本章局部生效。与 \cref{ch:planning} 的同字母（$\mu$ 测度、$\lambda$ 渗流强度）无关。本章用 $S_k$（与路径积分文献一致）记 rollout 代价，\cref{sec:mppi-family} 与 CEM 文献对齐处等价地写 $J_k$。
\end{note}

% ════════════════════════════════════════════════════════════════
\section{指数变换与 HJB 的线性化}
\label{sec:mppi-hjb}

\paragraph{动机：把不可解的非线性 HJB 整体线性化。}
最优控制的“主方程”是 HJB：理论上只要解出它，就得到全局最优反馈律。但它是一个非线性偏微分方程（PDE），状态维度 $n$ 一高，网格离散的格点数随 $n$ 指数爆炸（维度灾难），数值求解彻底不可行\cite{kappen2005linear}。梯度式 MPC（iLQR/DDP）退而求其次：不求全局值函数 $J$，只在当前轨迹附近做二次近似、反复迭代，代价是只能拿到\emph{局部}最优、且要求动力学可微。本节要走一条完全不同的路：对一大类\emph{控制仿射 + 二次控制代价}的随机问题，HJB 可经一次指数变换\emph{整体}线性化，而线性 PDE 有一个杀手锏——它的解可写成期望，能用蒙特卡洛估计。这是路径积分控制的起点，也是 MPPI 最终能“撒一批轨迹做加权平均”的根本原因。

\paragraph{反面：硬解非线性 HJB 是“宇宙年龄也算不完”。}
设想直接数值解非线性 HJB。即便只取 $n=12$（一个四旋翼的状态维度），每维粗糙地离散成 $50$ 个格点，状态空间就有 $50^{12}\approx2.4\times10^{20}$ 个格点；每个格点上还要对控制 $u$ 求极小、对时间反向积分。这不是“算得慢”，是不可能。历史上经典最优控制的精确解基本停在 LQR，正因为离开线性-二次特例，HJB 的非线性项就让一切解析与数值手段失效。

\paragraph{历史：Cole--Hopf 从 Burgers 方程到控制论。}
Kappen 的灵感来自物理\cite{kappen2005linear}。把非线性 HJB 变成线性方程的变换，与把经典 Hamilton--Jacobi 方程联系到薛定谔方程所用的变换是同一类：量子力学里 $\psi=e^{\mathrm{i}S/\hbar}$ 把作用量 $S$ 满足的非线性 HJ 方程变成线性薛定谔方程；最优控制里 $\psi=e^{-J/\lambda}$ 把值函数满足的非线性 HJB 变成线性 backward 扩散方程，温度 $\lambda$ 扮演 $\hbar$ 的角色。这种“作用量的指数”在数学上叫 Cole--Hopf 变换\cite{hopf1950partial,cole1951quasi}，最早用于把非线性 Burgers 方程线性化为热方程。这条思想链站在三代人肩上：源头是 Feynman 1948 的路径积分\cite{feynman1948spacetime}（“对所有路径求和、每条加权 $e^{\mathrm{i}S/\hbar}$”），Kac 1949 把它搬到扩散过程、证明抛物 PDE 的解可写成布朗运动路径上的期望（Feynman--Kac，\cref{sec:mppi-fk}）\cite{kac1949distributions}，Cole 与 Hopf 各自发现指数变换能线性化 Burgers 方程\cite{cole1951quasi,hopf1950partial}。Kappen 2005 看穿 HJB 恰好能套用这同一套工具\cite{kappen2005linear,kappen2005path}；Todorov 在离散侧独立给出镜像（LMDP，\cref{sec:mppi-lmdp}）\cite{todorov2009efficient}。

\subsection{从 HJB 到线性 backward PDE}
\label{sec:mppi-derivation}

我们一步步把非线性逼出来、再消掉，循序渐进地展开 Kappen 的推导。

\begin{definition}[控制仿射随机最优控制问题]\label{def:mppi-problem}
考虑控制仿射的随机微分方程（SDE）
\begin{equation}
  \mathrm{d}x=\big(f(x)+G(x)\,u\big)\,\mathrm{d}t+G(x)\,\mathrm{d}\boldsymbol{\xi},\qquad \mathrm{d}\boldsymbol{\xi}\sim\mathcal{N}(\mathbf{0},\ \boldsymbol{\Sigma}\,\mathrm{d}t),
  \label{eq:mppi-sde}
\end{equation}
其中 $f(x)$ 是\emph{漂移}（系统自身演化趋势），$G(x)$ 是\emph{控制矩阵}——\textbf{控制与噪声经同一通道 $G$ 注入}（这是“控制仿射”的含义，也是后面 Kappen 条件成立的前提）。代价取
\begin{equation}
  J=\mathbb{E}\!\left[\ \phi(x_T)+\int_0^T\Big(V(x,t)+\tfrac12\,u^\top R\,u\Big)\,\mathrm{d}t\ \right],
  \label{eq:mppi-cost}
\end{equation}
$\phi$ 是终端代价，$V(x,t)$ 是与状态有关的\emph{运行代价}（可任意复杂、非凸），$\tfrac12 u^\top R u$ 是\emph{二次控制代价}，$R\succ0$。\textbf{状态项 $V$ 与控制项分开}——状态可随便罚，控制必须二次罚。
\end{definition}

\noindent 这个结构是整个推导能走通的关键，后面反复用到。对最优值函数 $J(x,t)$（最优总代价就是初始时刻的 cost-to-go），动态规划给出随机 HJB 方程
\begin{equation}
  -\partial_t J=\min_u\Big\{\,V+\tfrac12 u^\top R u+(f+Gu)^\top\nabla_x J+\tfrac12\,\mathrm{tr}\!\big(G\boldsymbol{\Sigma}G^\top\,\nabla_x^2 J\big)\,\Big\}.
  \label{eq:mppi-hjb}
\end{equation}
最后一项 $\tfrac12\mathrm{tr}(G\boldsymbol{\Sigma}G^\top\nabla^2 J)$ 是\emph{随机性带来的二阶修正}（Itô 引理的产物）——确定性 HJB 没有、随机 HJB 独有；记住它，后面所有“魔法”都发生在它身上。

\begin{note}
\textbf{Itô vs Stratonovich.}\;那个 $\tfrac12$ 二阶项依赖于把 \cref{eq:mppi-sde} 按 Itô 约定解释。随机积分有两种主流约定：Itô 对 $\int g\,\mathrm{d}W$ 取黎曼和时用区间\emph{左端点}、链式法则多一个二阶项（即上面那项）、积分是鞅、便于概率推导；Stratonovich 用\emph{中点}、保留普通链式法则但漂移项要做对应修正。两者描述同一物理过程，只是记账方式不同。本章及绝大多数路径积分控制文献用 Itô。文献里若看到没有 $\tfrac12$ 项、却用普通链式法则的推导，多半是 Stratonovich，结论不矛盾。
\end{note}

\begin{theorem}[最优控制正比于值函数负梯度]\label{thm:mppi-optu}
\cref{eq:mppi-hjb} 大括号里对 $u$ 是标准二次型，求导置零 $Ru+G^\top\nabla J=0$ 得
\begin{equation}
  u^\star=-R^{-1}G^\top\nabla_x J.
  \label{eq:mppi-ustar}
\end{equation}
\end{theorem}
\begin{proof}
对 $\tfrac12 u^\top R u+u^\top G^\top\nabla J$ 求 $\partial/\partial u=Ru+G^\top\nabla J$，因 $R\succ0$ 该二次型严格凸，唯一极小在导数零点取得，解出 \cref{eq:mppi-ustar}。
\end{proof}
\noindent \cref{eq:mppi-ustar} 说明最优控制沿“代价下降最快”方向施加，强度由 $R^{-1}$（控制越便宜推得越狠）调制；这一步是确定性的，与随机性无关。

\paragraph{暴露非线性的根源。}把 $u^\star$ 代回 \cref{eq:mppi-hjb}。二次项 $\tfrac12(u^\star)^\top R u^\star+(u^\star)^\top G^\top\nabla J$，代入 $u^\star=-R^{-1}G^\top\nabla J$ 得 $\tfrac12\nabla J^\top GR^{-1}G^\top\nabla J-\nabla J^\top GR^{-1}G^\top\nabla J=-\tfrac12\nabla J^\top GR^{-1}G^\top\nabla J$，于是
\begin{equation}
  -\partial_t J=V+f^\top\nabla J\;\underbrace{-\;\tfrac12\,(\nabla J)^\top GR^{-1}G^\top\nabla J}_{\text{非线性项}}\;+\;\tfrac12\,\mathrm{tr}\!\big(G\boldsymbol{\Sigma}G^\top\nabla^2 J\big).
  \label{eq:mppi-hjb2}
\end{equation}
\textbf{全部麻烦集中在被标注的项上}：它对 $\nabla J$ 是二次的，使整个方程非线性。若能把它消掉，方程就线性了。这一步是诊断——我们已精确知道“敌人”是谁。

\begin{theorem}[Kappen 指数变换与线性化条件]\label{thm:mppi-kappen}
令 desirability $\psi(x,t)=\exp(-J(x,t)/\lambda)$，等价地 $J=-\lambda\log\psi$。当且仅当噪声协方差与控制代价的逆成正比、比例系数为温度 $\lambda$，
\begin{equation}
  \boxed{\ \boldsymbol{\Sigma}=\lambda R^{-1}\ }\qquad(\text{Kappen 条件}),
  \label{eq:mppi-kappen}
\end{equation}
\cref{eq:mppi-hjb2} 中两处“对 $\nabla\psi$ 二次”的项精确相消，HJB 化为关于 $\psi$ 的\emph{线性} backward PDE
\begin{equation}
  -\partial_t\psi=-\frac{V}{\lambda}\,\psi+f^\top\nabla\psi+\tfrac12\,\mathrm{tr}\!\big(G\boldsymbol{\Sigma}G^\top\,\nabla^2\psi\big),\qquad\psi(x,T)=\exp(-\phi(x)/\lambda).
  \label{eq:mppi-linearpde}
\end{equation}
\end{theorem}
\begin{proof}
为什么换到 $\psi$ 而非直接对 $J$ 下手？因为非线性恰好藏在对 $\nabla J$ \emph{二次}的那一项里。指数变换让“对 $\nabla J$ 二次”在求导后\emph{自动生出一个反向二次项}去抵消它。按链式法则（纯微积分）算三个导数：
\begin{equation}
  \partial_t J=-\lambda\frac{\partial_t\psi}{\psi},\quad
  \nabla J=-\lambda\frac{\nabla\psi}{\psi},\quad
  \nabla^2 J=-\lambda\Big(\frac{\nabla^2\psi}{\psi}-\frac{\nabla\psi\,\nabla\psi^\top}{\psi^2}\Big).
  \label{eq:mppi-chainrule}
\end{equation}
第三式多出的 $+\lambda\,\nabla\psi\nabla\psi^\top/\psi^2$ 也对 $\nabla\psi$ 二次，正好拿去与非线性项相消。代入 \cref{eq:mppi-hjb2}，两处二次项是
\[
  \underbrace{-\tfrac12\lambda^2\,\frac{\nabla\psi^\top GR^{-1}G^\top\nabla\psi}{\psi^2}}_{\text{来自非线性项}}
  \;+\;
  \underbrace{+\tfrac{\lambda}{2}\,\frac{\nabla\psi^\top G\boldsymbol{\Sigma}G^\top\nabla\psi}{\psi^2}}_{\text{来自 Hessian 项}}.
\]
两项相消当且仅当 $\lambda^2 GR^{-1}G^\top=\lambda\,G\boldsymbol{\Sigma}G^\top$，即 \cref{eq:mppi-kappen}。相消后乘开整理（$G\boldsymbol{\Sigma}G^\top=\lambda GR^{-1}G^\top$ 可回代）即得 \cref{eq:mppi-linearpde}；$\psi$、$\nabla\psi$、$\nabla^2\psi$ 都只出现一次幂，无任何乘积项，方程线性。终端条件由 $\psi(x,T)=e^{-J(x,T)/\lambda}=e^{-\phi(x)/\lambda}$ 得到。
\end{proof}

\noindent \textbf{为什么是 “backward” PDE？}\cref{eq:mppi-linearpde} 配的是\emph{终端}条件 $\psi(x,T)=e^{-\phi(x)/\lambda}$，而非初始条件——它从未来（$T$）向现在（$t$）反向求解。这不是技术怪癖，而是最优控制的本性：我们\emph{知道想在哪里结束}（终端代价 $\phi$ 刻画理想终点），却要倒推现在该怎么做。动态规划的 backward sweep（从终点反向递推值函数，\cref{thm:plan-bellman}）是同一逻辑的离散版本。$\psi$ 满足 backward 方程，正是 \cref{sec:mppi-fk} 能把它写成“从 $t$ 出发、向 $T$ 演化的轨迹期望”的根本原因。

\begin{insight}[指数变换的魔力：把 $\min$ 换成期望]
指数变换的本质，是\textbf{把“取极小（$\min$）”换成了“求期望（积分）”}。\cref{eq:mppi-hjb} 里有一个 $\min_u$，是非线性的源头之一；指数加权 $e^{-J/\lambda}$ 是 $\log\!\sum\!\exp$（软最小值，soft-min）的连续版本——当 $\lambda\to0$ 时 $-\lambda\log\int e^{-J/\lambda}\to\min J$，精确退回硬最小值。换句话说，我们没有“解决”非线性，而是把它\emph{搬进了指数的内部}，让外层运算（积分、期望）变成线性的。这个“用 log-sum-exp 把 $\min$ 软化、从而线性可加”的思想，会在 \cref{sec:mppi-freeenergy} 的自由能对偶里以完全相同的面貌再现——届时你会发现 Kappen 的 PDE 推导与 Williams 的信息论推导是同一枚硬币的两面。
\end{insight}

\subsection{Kappen 条件的物理意义}
\label{sec:mppi-kappen-physics}

$\boldsymbol{\Sigma}=\lambda R^{-1}$ 不是可有可无的技术细节，它编码一句很强的物理陈述：\textbf{在噪声大的方向上，控制是便宜的；在噪声小的方向上，控制是昂贵的。}因为 $\boldsymbol{\Sigma}$ 大对应 $R$ 小（控制便宜）、$\boldsymbol{\Sigma}$ 小对应 $R$ 大（控制昂贵），二者此消彼长。直觉上很合理——若某控制通道天生被大量噪声搅动，再花大力气精确控制它也徒劳，系统“乐于”让你在那个方向上随便推；反过来，几乎无噪声的通道每一点控制都精确生效，理应谨慎使用。这个条件也解释了为何 \cref{sec:mppi-fk} 之后能用“无控系统的随机轨迹”估计最优控制：噪声的分布与控制的代价是\emph{同一个几何对象}（差一个 $\lambda$），所以“系统在噪声驱动下自发探索的方向”恰好就是“控制代价低、值得探索的方向”。

\begin{table}[H]\centering\small
\caption{一族横跨三领域的同构：经典力学、量子力学与随机最优控制}
\label{tab:mppi-trinity}
\begin{tabular}{p{0.2\textwidth}p{0.24\textwidth}p{0.22\textwidth}p{0.24\textwidth}}
\toprule
 & 经典力学（HJ） & 量子力学（薛定谔） & 随机最优控制（本章） \\
\midrule
核心未知量 & 作用量 $S$ & 波函数 $\psi$ & 值函数 $J$ / desirability $\psi=e^{-J/\lambda}$ \\
控制方程 & HJ（非线性一阶 PDE） & 薛定谔（线性） & HJB（非线性）$\to$ 线性 backward PDE \\
线性化的桥 & —— & $\psi=e^{\mathrm{i}S/\hbar}$（WKB） & $\psi=e^{-J/\lambda}$（Cole--Hopf） \\
“小参数” & —— & 普朗克常数 $\hbar$ & 温度 $\lambda$ \\
参数 $\to0$ 极限 & —— & $\hbar\to0$：回到经典力学 & $\lambda\to0$：回到确定性最优控制（PMP） \\
求解手段 & 特征线 / 最小作用量 & 路径积分 & 路径积分（对被动轨迹采样，\cref{sec:mppi-fk}） \\
\bottomrule
\end{tabular}
\end{table}

\noindent 最后一行值得停一下。量子力学里 $\hbar\to0$ 回到经典确定性世界；最优控制里 $\lambda\to0$（噪声趋零）回到确定性最优控制——Kappen 证明，此极限下路径积分的 Laplace 近似由\emph{作用量极小的那条路径}主导、且在 $\nu\to0$ 时\emph{精确}，给出的边值 Hamilton 方程正是 Pontryagin 极小值原理（PMP）的常微分方程，确定性最优控制律由此恢复\cite{kappen2005path,kappen2005linear}。你在确定性最优控制课里学过的 PMP，正是本章随机理论在零温极限下的特例；噪声一旦打开，PMP 的单条最优轨迹就“软化”成一个轨迹\emph{分布}——而这分布，正是 \cref{sec:mppi-fk} 起要采样的对象。

\begin{pitfall}[概念误区：以为指数变换对任意最优控制问题都能线性化]
看到 $\psi=e^{-J/\lambda}$ 把 HJB 变线性，就以为这是个万能技巧、对任何代价任何动力学都成立。\textbf{后果}：在控制代价非二次、或噪声不进控制通道的问题上套用，得到的“线性 PDE”其实残留未消去的非线性项，后续所有结论（Feynman--Kac、蒙特卡洛估计）全部失效，却很难一眼看出错在哪。\textbf{根本原因}：线性化严格依赖 \cref{thm:mppi-kappen} 的相消，而相消\emph{只在 Kappen 条件 $\boldsymbol{\Sigma}=\lambda R^{-1}$ 下发生}，这要求 (a) 控制仿射、(b) 控制代价二次、(c) 噪声与控制走同一通道 $G$，三者缺一不可。\textbf{对策}：套用前核对这三个前提；若控制代价非二次或动力学非控制仿射，应改用 \cref{sec:mppi-freeenergy} 的信息论推导——Williams 2018 正为去掉这些限制而生\cite{williams2018information}。
\end{pitfall}

\begin{pitfall}[思维陷阱：把 Kappen 条件当“可事后调”的超参数]
以为 $\boldsymbol{\Sigma}$、$R$、$\lambda$ 可各自独立设定、跑不好再分别微调。\textbf{后果}：把噪声方差、控制权重、温度当三个旋钮乱拧，结果要么数值发散，要么求出的控制律根本不是任何问题的最优解。\textbf{根本原因}：$\boldsymbol{\Sigma}=\lambda R^{-1}$ 是线性化成立的\emph{充要条件}，不是建议；一旦三者不满足这条等式，你求出的 $\psi$ 就不再是原问题的解。\textbf{对策}：把其中两个当自由量、第三个由等式定出——常见做法是固定 $\boldsymbol{\Sigma}$ 与 $R$、令 $\lambda$ 由问题尺度决定；或在 MPPI 实现里把控制代价吸收进 $\lambda$，只保留一个温度旋钮。
\end{pitfall}

\begin{practice}
\begin{enumerate}
  \item 补全 \cref{thm:mppi-kappen} 证明省略的代数：把 \cref{eq:mppi-chainrule} 三个导数完整代入 \cref{eq:mppi-hjb2}，逐项写出，验证除两个“对 $\nabla\psi$ 二次”的项外，其余整理后恰好给出 \cref{eq:mppi-linearpde}，并标出每项来源。
  \item 证明确定性极限：$\boldsymbol{\Sigma}\to\mathbf{0}$（于是 $\lambda\to0$）时路径积分解经 Laplace 近似退化为 PMP。提示：$-\lambda\log\int e^{-J/\lambda}\,\mathrm{d}\tau\xrightarrow{\lambda\to0}\min_\tau J$。
  \item （开放）设 Kappen 条件被破坏：$\boldsymbol{\Sigma}=c\,\lambda R^{-1}$ 但 $c\neq1$。\cref{thm:mppi-kappen} 的两项不再相消，残留一个怎样的项？它对 $\psi$ 还线性吗？这对“用无控轨迹做蒙特卡洛估计”意味着什么困难？（定性即可——这正是 Williams 2018 要绕开的麻烦。）
\end{enumerate}
\end{practice}

% ════════════════════════════════════════════════════════════════
\section{Feynman--Kac 与路径积分控制律}
\label{sec:mppi-fk}

\paragraph{动机：根本不解 PDE，而是采样求平均。}
\cref{sec:mppi-hjb} 把方程变简单了，但线性 PDE 仍是 PDE——高维下直接数值解它，维度灾难一分没少（$n=12$ 时无论线不线性都要 $50^{12}$ 个格点）。我们需要的不是“更好解的 PDE”，而是\emph{根本不解 PDE}。Feynman--Kac 引理正是这把钥匙：它在数学上建立一座桥——\textbf{任何线性抛物 backward PDE 的解，都可写成某个随机过程上的期望}\cite{kac1949distributions}。一旦写成期望，求解方式就从“在整个状态空间铺网格、反向扫描”变成“从当前点出发、向前撒一批随机轨迹、取平均”：前者计算量随维度指数爆炸，后者只取决于采样条数 $K$、\emph{与状态维度无关}（蒙特卡洛的维度无关性）。这一步是路径积分控制全部实用价值的来源。

\begin{theorem}[Feynman--Kac 路径积分公式]\label{thm:mppi-feynmankac}
\cref{eq:mppi-linearpde} 可写成 $-\partial_t\psi=(\mathcal{L}-\tfrac{V}{\lambda})\psi$，其中 $\mathcal{L}=f^\top\nabla+\tfrac12\mathrm{tr}(G\boldsymbol{\Sigma}G^\top\nabla^2)$ 正是\emph{无控扩散过程}（被动动力学，passive dynamics）
\begin{equation}
  \mathrm{d}x=f(x)\,\mathrm{d}t+G(x)\,\mathrm{d}\boldsymbol{\xi}
  \label{eq:mppi-passive}
\end{equation}
的无穷小生成元。其解由对被动轨迹的期望给出：
\begin{equation}
  \boxed{\;\psi(x,t)=\mathbb{E}_{\mathbb{Q}}\!\left[\ \exp\!\Big(-\frac{\phi(x_T)}{\lambda}-\int_t^T\frac{V(x_s,s)}{\lambda}\,\mathrm{d}s\Big)\ \Bigg|\ x_t=x\ \right]\;}
  \label{eq:mppi-fk}
\end{equation}
其中 $\mathbb{Q}$ 是 \cref{eq:mppi-passive} 的轨迹分布，期望对从 $x_t=x$ 出发的所有\emph{无控}随机轨迹取平均。
\end{theorem}

\begin{derivation}[为什么 PDE 的解是期望：离散随机游走视角]
\label{der:mppi-fk}
不需测度论，用离散随机游走就能看清这个“PDE $\leftrightarrow$ 期望”的对应从何而来。把时间切成小步 $\Delta t$，无控系统一步内从 $x$ 随机跳到邻近的 $x'$（高斯转移）。\cref{eq:mppi-linearpde} 离散化后是一条递推：当前格点的 desirability 等于“下一步各落点 desirability 的期望”，再乘这一步的局部代价因子 $e^{-V(x)\Delta t/\lambda}$：
\begin{equation}
  \psi(x,t)\ \approx\ e^{-V(x)\Delta t/\lambda}\ \mathbb{E}_{x'\sim p(\cdot\mid x)}\big[\psi(x',t+\Delta t)\big].
  \label{eq:mppi-fk-discrete}
\end{equation}
零阶项 $-V\psi/\lambda$ 变成乘性因子，扩散项变成对随机落点求期望。把这条递推从 $t$ 一路展开到终端 $T$：每多走一步就多乘一个 $e^{-V\Delta t/\lambda}$、多套一层期望，$T$ 处用终端条件 $\psi(x_T,T)=e^{-\phi(x_T)/\lambda}$ 收尾。所有一步因子连乘起来，正是 $\exp(-\frac1\lambda[\phi(x_T)+\sum_t V\Delta t])$；所有嵌套期望合起来，正是“对整条随机游走轨迹求期望”。$\Delta t\to0$ 取极限，求和变积分、随机游走变布朗运动，就回到 \cref{eq:mppi-fk}。\textbf{一句话}：backward PDE 的解之所以是期望，是因为解 PDE 等价于把“逐格反向递推”沿所有随机路径展开成连乘，而沿路径连乘再对路径平均，就是路径积分。
\end{derivation}

\noindent 从这里开始，所有“采样”都是对\emph{无控系统} \cref{eq:mppi-passive} 采样，而非对受控系统采样——记住这一点。\cref{eq:mppi-fk} 的物理含义：desirability $\psi(x,t)$ 等于“从 $x$ 出发、不施加任何控制、任系统随噪声漂流、所积累代价的指数的平均”。代价低的轨迹贡献大的 $e^{-S/\lambda}$，代价高的趋于零——所以 $\psi$ 大意味着“从这里随便走走就能撞上低代价的未来”，这个状态天生合意。\textbf{最优代价被编码成了被动轨迹的一个加权统计量。}

\begin{theorem}[路径积分控制律]\label{thm:mppi-control-law}
由 \cref{eq:mppi-ustar} 与 $J=-\lambda\log\psi$（于是 $\nabla J=-\lambda\nabla\psi/\psi$），最优控制为
\begin{equation}
  u^\star(x,t)=-R^{-1}G^\top\nabla J=\lambda R^{-1}G^\top\,\frac{\nabla\psi}{\psi}=\lambda R^{-1}G^\top\,\nabla\log\psi.
  \label{eq:mppi-pi-control}
\end{equation}
把 \cref{eq:mppi-fk} 代入 $\nabla\psi/\psi$，最优控制变成\emph{对无控轨迹的一个加权平均}：权重正比于每条轨迹的 $e^{-S/\lambda}$，被加权的量是该轨迹初始时刻所需的“推力”。
\end{theorem}

\noindent 这就是路径积分控制律——\textbf{最优控制不是解出来的，是采样平均出来的}。MPPI 的全部，就是把这个连续时间的期望，换成离散时间下 $K$ 条 rollout 的有限和。

\subsection{贯穿例子：一维点质量（路径积分 = LQR 的硬证据）}
\label{sec:mppi-pointmass}

抽象公式要落到一个能亲手算的例子上才站得住。取最简单的控制仿射系统——一维点质量 / 单积分器
\begin{equation}
  \mathrm{d}x=u\,\mathrm{d}t+\sigma\,\mathrm{d}W,\qquad f=0,\ G=1,\ \boldsymbol{\Sigma}=\sigma^2,
  \label{eq:mppi-pm-sde}
\end{equation}
即“位置 $x$ 的变化率就是控制 $u$，外加强度 $\sigma$ 的白噪声”。控制代价 $\tfrac12 r u^2$（$R=r$），运行代价 $V=0$，终端代价 $\phi(x)=\tfrac{a}{2}x^2$（想把粒子拉回原点）。Kappen 条件在此即 $\sigma^2=\lambda/r$，即 $\lambda=r\sigma^2$。无控动力学是纯布朗运动 $\mathrm{d}x=\sigma\,\mathrm{d}W$，从 $x_t=x$ 出发的终点服从 $x_T\mid x_t=x\sim\mathcal{N}(x,\sigma^2(T-t))$。代入 \cref{eq:mppi-fk}（$V=0$，只剩终端项）得高斯积分，用恒等式 $\mathbb{E}_{y\sim\mathcal{N}(\mu,s^2)}[e^{-by^2/2}]=(1+bs^2)^{-1/2}\exp(-\tfrac{b\mu^2}{2(1+bs^2)})$（取 $b=a/\lambda,\mu=x,s^2=\sigma^2(T-t)$）：
\begin{equation}
  \psi(x,t)=\frac{1}{\sqrt{1+\tfrac{a}{\lambda}\sigma^2(T-t)}}\ \exp\!\left(-\frac{(a/\lambda)\,x^2}{2\big(1+\tfrac{a}{\lambda}\sigma^2(T-t)\big)}\right).
  \label{eq:mppi-pm-psi}
\end{equation}
对 $\log\psi$ 求梯度，再乘 $\lambda R^{-1}G^\top=\lambda/r=\sigma^2$，得最优控制律
\begin{equation}
  u^\star(x,t)=\sigma^2\,\partial_x\log\psi=-\,\frac{(a/r)\,x}{1+(a/r)(T-t)}=-k(t)\,x,
  \label{eq:mppi-pm-ustar}
\end{equation}
一个\emph{时变线性反馈}，增益 $k(t)=\tfrac{a/r}{1+(a/r)(T-t)}$ 随 $t\to T$（临近终点）单调增大——越接近终点纠偏越狠，符合直觉。

\begin{theorem}[路径积分在 LQ 情形等于 Riccati 解]\label{thm:mppi-lqr}
\cref{eq:mppi-pm-sde} 恰是一个线性二次（LQR）问题（$A=0,B=1$，状态代价 $0$，控制代价 $r$，终端代价 $a$）。其连续时间 Riccati 方程 $\dot P=P^2/r,\ P(T)=a$ 解得 $P(t)=\tfrac{ar}{a(T-t)+r}$，对应最优反馈 $u^\star_{\mathrm{LQR}}=-\tfrac1r P(t)x=-\tfrac{a/r}{1+(a/r)(T-t)}x$，与 \cref{eq:mppi-pm-ustar} \emph{逐字相同}。
\end{theorem}

\noindent 路径积分——一个表面上完全不碰 Riccati、只靠“对布朗轨迹取期望”的方法——在 LQ 情形给出和经典最优控制\emph{一模一样}的答案。这不是巧合：LQR 是路径积分能精确闭式求解的特例。而路径积分的真正威力在于，当 $\phi$ 或 $V$ \emph{不是}二次的（比如 $\phi$ 是“撞墙处跳到 $\infty$”的不连续障碍代价）、Riccati 彻底失效时，那个高斯积分换成蒙特卡洛求和\emph{依然成立}。下面这段最朴素的蒙特卡洛直接对被动布朗轨迹采样、按 desirability 加权，不解任何方程、不碰 Riccati：

\begin{codebox}[Python]{路径积分 = LQR 的蒙特卡洛验证（一维点质量）}
import numpy as np
a, r, T, dt = 5.0, 1.0, 1.0, 0.02
sigma = 0.3; lam = r * sigma**2              # Kappen condition  lambda = r*sigma^2
H = int(round(T / dt))

def pi_control(x0, t_idx, K=400_000):        # path-integral estimate of u* at (t_idx, x0)
    steps = H - t_idx
    xT = x0 + (sigma*np.sqrt(dt)*np.random.randn(K, steps)).sum(axis=1)  # passive endpoints
    w = np.exp(-0.5*a*xT**2 / lam); w /= w.sum()         # desirability weights  ~ exp(-phi/lam)
    return -(a/r) * (w*xT).sum()             # Monte-Carlo estimate of  u* = lam d_x log psi

for t in [0.0, 0.5, 0.98]:                   # x0 = 1
    ti = int(round(t/dt)); k = (a/r)/(1 + (a/r)*(T-t))   # closed-form Riccati gain
    print(f"t={t:.2f}  path-integral {pi_control(1.0, ti):+.4f}   closed-form {-k:+.4f}")
\end{codebox}

\noindent 输出（$K=4\times10^5$，$x_0=1$）为 \texttt{t=0.00: $-0.8455$ /\allowbreak  $-0.8333$}、\texttt{t=0.50: $-1.3853$ /\allowbreak  $-1.4286$}、\texttt{t=0.98: $-4.5407$ /\allowbreak  $-4.5455$}：三个时刻蒙特卡洛估计都贴着闭式 Riccati 解，残差来自有限采样的方差，$K$ 调大还会继续缩小。它真正的意义在于：把终端代价换成任意（哪怕不连续的）代价，这段代码\emph{一字不改}照样跑，可 Riccati 在那时早已失效。

\subsection{从期望到蒙特卡洛估计器，及方差爆炸的隐患}
\label{sec:mppi-mc}

闭式高斯积分只在 LQ 特例存在。一般情况下退回 \cref{eq:mppi-fk} 的定义，用 $K$ 条无控轨迹的样本均值估计期望：
\begin{equation}
  \hat\psi(x,t)=\frac{1}{K}\sum_{k=1}^{K}\exp\!\Big(-\frac{\phi(x_T^{(k)})}{\lambda}-\int_t^T\frac{V(x_s^{(k)})}{\lambda}\,\mathrm{d}s\Big),
  \label{eq:mppi-psihat}
\end{equation}
每条 $x^{(k)}$ 是一条独立模拟的无控轨迹；由大数定律 $K\to\infty$ 时 $\hat\psi\to\psi$。\textbf{这里出现了 MPPI 的雏形}：撒轨迹、按 $e^{-\text{代价}/\lambda}$ 加权、求平均。但纯被动采样有一个致命弱点，它直接催生 \cref{sec:mppi-freeenergy}。设想终端代价要求粒子精确停在某个远离原点的目标点：无控的布朗运动\emph{几乎不可能}自己漂到那里，于是绝大多数轨迹 $e^{-S/\lambda}\approx0$，只有极少数侥幸轨迹贡献几乎全部权重——估计方差爆炸，$K$ 要大到天文数字才够。解决办法是\emph{别再从无控分布 $\mathbb{Q}$ 采样，而从一个已经接近最优的分布采样}（比如“当前控制序列 + 噪声”），再用重要性权重修正。这一步叫 generalized importance sampling，是 Williams 2017 把 Kappen 理论变成实用算法的关键转折\cite{williams2017model}，也是 \cref{sec:mppi-freeenergy} 的主线。

\begin{insight}[backward sweep $\to$ forward sampling：MPPI 能用 GPU 的根源]
路径积分控制把动态规划的 \emph{backward sweep} 换成了 \emph{forward sampling}。动态规划（Riccati、值迭代）必须从终点出发、反向逐状态计算每个状态的最优值——本质是串行的、要遍历状态空间。路径积分反过来：从当前点出发、向前并行撒一批轨迹、取加权平均——本质是并行的、只看采到的轨迹。这个范式翻转不仅消解了维度灾难，更\emph{直接决定了 MPPI 的硬件形态}：$K$ 条 forward rollout 彼此独立，是 GPU 数千线程的完美负载。一句话——路径积分之所以能用 GPU，是因为 Feynman--Kac 把“解方程”变成了“撒点”。
\end{insight}

\begin{pitfall}[概念误区：以为 Feynman--Kac 直接给出“最优轨迹”]
把 $\psi(x,t)$ 的期望表达式理解成“它算出了那条最优路径”。\textbf{后果}：去代码里找“最优轨迹变量”，或把采样得到的某条低代价 rollout 当最优解直接执行，得到次优甚至危险的控制。\textbf{根本原因}：Feynman--Kac 给的是一个\emph{标量场} $\psi(x,t)$（每个状态的合意度），最优控制由它的对数梯度 $\nabla\log\psi$ 给出，是一个\emph{反馈律}，不是任何单条轨迹；采样到的轨迹只是估计这个场/梯度的“原料”。\textbf{对策}：始终区分“采样轨迹（手段）”与“加权平均得到的控制（结果）”——MPPI 执行的是 $K$ 条轨迹\emph{加权融合}后的控制序列，绝不是某一条 rollout。
\end{pitfall}

\begin{practice}
\begin{enumerate}
  \item 对一维点质量 + 二次终端代价，独立完成 \cref{sec:mppi-pointmass} 省略的高斯积分得 \cref{eq:mppi-pm-psi}，再求 $u^\star$；独立解一遍标量 Riccati，验证反馈增益\emph{逐项相等}。说明：$\phi$ 非二次时还会相等吗？为什么？
  \item 直接验证 \cref{eq:mppi-fk} 确实满足 \cref{eq:mppi-linearpde}：把 $\psi=\mathbb{E}_{\mathbb{Q}}[\cdot]$ 对 $t$ 求导，用 Itô 引理处理期望内的随机积分，证明它满足 $-\partial_t\psi=-\tfrac{V}{\lambda}\psi+\mathcal{L}\psi$ 且满足终端条件。提示：把 $e^{-\int_t^T V/\lambda\,\mathrm{d}s}\psi$ 视为鞅。
  \item （开放）若不从无控分布、而是给系统施加标称控制 $\bar u(t)$ 后再采样（从“有漂移”的分布采样），Feynman--Kac 的期望要怎样修正才能仍估计同一个 $\psi$？该乘一个依赖 $\bar u$ 哪些量的“修正因子”？（定性即可——这就是 \cref{sec:mppi-freeenergy} 重要性权重的前身，严格形式是 Girsanov 变换。）
\end{enumerate}
\end{practice}

% ════════════════════════════════════════════════════════════════
\section{离散时间的镜像：线性可解 MDP}
\label{sec:mppi-lmdp}

\cref{sec:mppi-hjb,sec:mppi-fk} 全程在连续时间里展开（SDE、HJB PDE、Feynman--Kac）。但同一个“指数变换 $\to$ 线性化”的思想，在\emph{离散时间的 MDP} 上有一个同样漂亮、更贴近 RL 背景的版本——Todorov 2009 的线性可解 MDP（Linearly-Solvable MDP, LMDP）\cite{todorov2009efficient}。它不走 PDE 路线，而是直接改造贝尔曼方程；理解它有两重价值：补全路径积分的离散侧，并为 \cref{sec:mppi-freeenergy} 的自由能--KL 对偶埋下伏笔。

\paragraph{反面：标准 MDP 的 $\min_u$ 不可解析。}标准 MDP 的贝尔曼方程 $v(x)=\min_u\{c(x,u)+\mathbb{E}_{x'}[v(x')]\}$ 里的 $\min_u$ 是离散、组合式的——动作是“向左/向右”这类符号，对它求极小只能穷举，无解析解。Todorov 的问题与 Kappen 如出一辙：能不能像 LQR 那样，\emph{构造一类 MDP 让 $\min_u$ 有闭式解}？

\begin{definition}[线性可解 MDP]\label{def:mppi-lmdp}
LMDP 做两件事。其一，把动作从“符号”换成“对下一状态的分布”——控制直接是受控转移概率 $u(x'\mid x)$。其二，引入被动动力学 $p(x'\mid x)$（系统无控时的自然转移），并把控制代价\emph{规定}为受控转移与被动转移之间的 KL 散度：
\begin{equation}
  c(x,u)=q(x)+\lambda\,\mathrm{KL}\big(u(\cdot\mid x)\,\big\|\,p(\cdot\mid x)\big),
  \label{eq:mppi-lmdp-cost}
\end{equation}
$q(x)$ 是任意状态代价（不受限制），控制代价则度量“你把系统推离它自然行为有多远”——KL 越大推得越狠越贵，且 $\mathrm{KL}\ge0$、等号当且仅当 $u=p$（完全不控制），符合直觉。
\end{definition}

\begin{theorem}[线性贝尔曼方程 $z=MPz$]\label{thm:mppi-lmdp}
代入 \cref{eq:mppi-lmdp-cost}，贝尔曼方程成为 $v(x)=q(x)+\min_{u(\cdot\mid x)}\{\lambda\,\mathrm{KL}(u\|p)+\mathbb{E}_{x'\sim u}[v(x')]\}$。大括号里是一个对分布 $u$ 的 KL 正则最小化，其最优解是吉布斯型的——把被动转移按 $e^{-v(x')/\lambda}$ 重新加权。定义 desirability $z(x)=e^{-v(x)/\lambda}$（与连续侧 $\psi=e^{-J/\lambda}$ 一一对应），最优受控转移有闭式
\begin{equation}
  u^\star(x'\mid x)=\frac{p(x'\mid x)\,z(x')}{\sum_{x''}p(x''\mid x)\,z(x'')}=\frac{p(x'\mid x)\,z(x')}{[Pz](x)},
  \label{eq:mppi-lmdp-ustar}
\end{equation}
即被动转移被 desirability 重新加权再归一化。把 $u^\star$ 代回贝尔曼方程整理得关于 $z$ 的\emph{线性}方程
\begin{equation}
  \boxed{\,z(x)=e^{-q(x)/\lambda}\,[Pz](x)\,}\qquad\Longleftrightarrow\qquad z=MPz,\quad M=\mathrm{diag}\big(e^{-q/\lambda}\big),
  \label{eq:mppi-zmpz}
\end{equation}
$P$ 是被动转移矩阵。
\end{theorem}

\noindent 非线性的贝尔曼方程被驯服成一个线性代数问题：首次到达（first-exit）问题是线性方程组，无穷时域折扣/平均代价问题是\emph{最大特征值问题}（$z$ 是 $MP$ 的 Perron 主特征向量，对应最大特征值 $\rho$ 给出最优平均每步代价 $-\lambda\log\rho$）。Todorov 由此得到求最短路径的 $O(n)$ 算法，并指出这是“第一类能在给定最优代价后解析求出最优动作的一般 MDP”\cite{todorov2009efficient}。

\begin{insight}[KL 控制代价 = 离散侧的 Kappen 条件]
LMDP 的全部魔力来自“把控制代价定义成 KL 散度”这个选择。连续时间里（\cref{sec:mppi-hjb}）是 Kappen 条件 $\boldsymbol{\Sigma}=\lambda R^{-1}$（二次控制代价 + 噪声进控制通道）让线性化发生；离散时间里，是 KL 形式的控制代价扮演同样的角色。两者并非巧合——KL 散度在小扰动下的二阶展开恰是一个马氏距离型的二次型，所以“二次控制代价”与“KL 控制代价”在局部是同一个东西。这预告 \cref{sec:mppi-freeenergy}：那里 KL 控制代价不必“人为规定”，它会从轨迹空间的自由能变分原理里\emph{自然涌现}——而一旦控制代价是 KL，二次代价的限制就被彻底甩掉（这正是 Williams 2018 推广 MPPI 的关键一步）。
\end{insight}

\begin{table}[H]\centering\small
\caption{连续与离散：同一枚硬币的两面}
\label{tab:mppi-cont-disc}
\begin{tabular}{lll}
\toprule
维度 & 连续时间（Kappen，\cref{sec:mppi-hjb,sec:mppi-fk}） & 离散时间（Todorov LMDP） \\
\midrule
让线性化发生的代价结构 & 二次控制代价 + Kappen 条件 $\boldsymbol{\Sigma}=\lambda R^{-1}$ & 控制代价 $=$ KL 散度 $\lambda\,\mathrm{KL}(u\|p)$ \\
指数变换 & $\psi=e^{-J/\lambda}$ & $z=e^{-v/\lambda}$ \\
线性化结果 & 线性 backward PDE & 线性贝尔曼 $z=MPz$ \\
求解手段 & Feynman--Kac $\to$ 蒙特卡洛采样 & 线性方程组 / 最大特征值 \\
“被动”基准 & 无控扩散 $\mathrm{d}x=f\mathrm{d}t+G\mathrm{d}\boldsymbol{\xi}$ & 被动转移 $p(x'\mid x)$ \\
最优控制 & $u^\star=\lambda R^{-1}G^\top\nabla\log\psi$ & $u^\star(x'\mid x)\propto p(x'\mid x)z(x')$ \\
\bottomrule
\end{tabular}
\end{table}

\noindent \textbf{两个额外的结构红利。}其一是\emph{可叠加性（compositionality）}：因为 \cref{eq:mppi-zmpz} 对 $z$ 线性，若干“基础任务”（仅终端代价不同）的最优 desirability 可\emph{线性加权}组合出复合任务的最优解——这在标准 MDP 里做不到（值函数满足带 $\min$ 的非线性贝尔曼方程，两任务最优值函数的线性组合一般不是复合任务的最优值），是 LMDP 独有的好处\cite{todorov2009compositionality}。其二是 \emph{Z-learning}：可像 Q-learning 那样从经验里 off-policy 地学 $z$，且因 $z$ 只是状态的函数（不需状态-动作值），Todorov 证明 Z-learning 比 Q-learning 更高效\cite{todorov2009efficient}。这两点都源自同一事实——\emph{线性}：一旦把非线性的 $\min$ 换成线性的指数加权，叠加原理就回来了。

\begin{practice}
\begin{enumerate}
  \item 补全 LMDP 线性化：把 \cref{eq:mppi-lmdp-ustar} 代回带 KL 控制代价的贝尔曼方程，证明它化简为 \cref{eq:mppi-zmpz}。提示：KL 正则最小化的最优解是吉布斯分布，其最优值等于 $-\lambda\log[Pz](x)$（配分函数的对数，又一次 log-sum-exp）。完成后对照连续推导，指出“KL 控制代价”在离散侧扮演“Kappen 条件”在连续侧的角色。
  \item 在一条 $5$ 状态链上（一端为吸收目标 $z=1,q=0$，其余每步罚 $q=1$，被动动力学为左右等概率随机游走）数值解 $z=MPz$（幂迭代），再由 $z$ 还原 $v=-\lambda\log z$ 与最优策略 \cref{eq:mppi-lmdp-ustar}，验证 $z$ 从目标向远端指数衰减、最优转移把概率压向更接近目标的状态。
  \item （叠加性）取一条 $7$ 状态链、两端可作目标。分别解“到左端”与“到右端”的 $z_A,z_B$，再解“到任意一端”（终端 desirability 取两者等权混合）的 $z_C$，验证 $z_C=\tfrac12 z_A+\tfrac12 z_B$ 逐点相等。说明这在标准 MDP 里为何不成立。
\end{enumerate}
\end{practice}

% ════════════════════════════════════════════════════════════════
\section{自由能--KL 对偶与 MPPI 权重公式}
\label{sec:mppi-freeenergy}

\paragraph{动机：抛开 HJB，对任意黑箱代价重得路径积分。}
\cref{sec:mppi-hjb,sec:mppi-fk} 的推导很美，但适用范围被三个前提卡死：控制仿射、二次控制代价、噪声进控制通道（Kappen 条件）。现代采样式 MPC 想做的恰恰是突破这些前提——动力学是一个 MuJoCo 黑箱或一张神经网络，代价里含“碰撞即 $+\infty$”的不连续 indicator，控制代价未必二次。这些场景下 Kappen 的 PDE 推导根本写不出来（HJB 都不一定存在经典解）。Williams 2018 的关键洞察是：\textbf{我们其实根本不需要 HJB}\cite{williams2018information}。路径积分控制的核心结论——“按 $e^{-\text{代价}/\lambda}$ 给轨迹加权”——可\emph{完全脱离 PDE}，仅从信息论的一个变分原理（自由能--KL 对偶）推出来。这个推导只用到“轨迹有个代价、我们能从某个分布采样轨迹”两件事，对动力学和代价的形式毫无要求。于是同一套加权机制瞬间从“一小类 LQ-like 问题”推广到“任何能采样、能算代价的问题”。

\paragraph{历史：从 PI$^2$ 到信息论 MPPI 的三代传承。}
Theodorou、Buchli、Schaal 2010 在 JMLR 提出 PI$^2$（Policy Improvement with Path Integrals）：把路径积分控制改造成策略搜索算法，用 reward-weighted regression 更新策略参数，\emph{没有学习率、没有矩阵求逆}——这种“按代价指数加权求平均”的更新，是 MPPI 的直接祖先\cite{theodorou2010generalized}。Williams、Aldrich、Theodorou 2017 在 JGCD 引入 generalized importance sampling，允许改变采样分布的漂移与扩散（从“当前控制序列 + 噪声”采样而非裸采无控轨迹），并用 GPU 并行，\emph{MPPI 算法正式诞生}\cite{williams2017model}。Williams 等 2018 在 IEEE T-RO 用自由能--KL 对偶给出集大成的信息论推导，去掉 control-affine 与二次代价限制，并在 AutoRally 上跑了 100+ 公里实车\cite{williams2018information}。谱系：Kappen（PDE 起源）$\to$ Theodorou（PI$^2$ 策略搜索）$\to$ Williams（importance sampling + GPU + 信息论统一）。

\subsection{轨迹空间上的自由能与对偶定理}
\label{sec:mppi-duality}

\paragraph{把问题搬到轨迹空间。}不再逐时刻看状态，而是把整条轨迹 $\tau=(x_0,x_1,\dots,x_T)$ 当成一个对象。给定一个\emph{基分布}（base / proposal distribution）$\mathbb{Q}$——在 MPPI 里就是“当前控制序列 $U=\{u_0,\dots,u_{T-1}\}$ 叠加高斯噪声 $\boldsymbol{\varepsilon}\sim\mathcal{N}(\mathbf{0},\boldsymbol{\Sigma})$ 后 rollout 出的轨迹分布”。每条轨迹有一个\emph{总代价}
\begin{equation}
  S(\tau)=\phi(x_T)+\sum_{t=0}^{T-1}V(x_t,t)\,\Delta t.
  \label{eq:mppi-S}
\end{equation}
注意：$S(\tau)$ 只是一个\emph{标量数值}——它怎么算出来（黑箱仿真器、神经网络、含不连续 indicator 的代价）\emph{完全无所谓}，后面所有推导都不对 $S$ 求导。这正是脱离 Kappen 假设的关键。

\begin{definition}[轨迹分布上的自由能]\label{def:mppi-freeenergy}
借统计物理，定义轨迹分布 $\mathbb{Q}$ 上的自由能（free energy）
\begin{equation}
  \mathcal{F}(\mathbb{Q})=-\lambda\,\log\,\mathbb{E}_{\mathbb{Q}}\!\big[\,e^{-S(\tau)/\lambda}\,\big].
  \label{eq:mppi-F}
\end{equation}
这就是统计物理里 $F=-k_B T\log Z$ 的形式（$\lambda$ 当温度，$\mathbb{E}_{\mathbb{Q}}[e^{-S/\lambda}]$ 当配分函数 $Z$），把“所有轨迹的代价”用温度 $\lambda$ 软聚合成单个标量：$\lambda\to0$ 时 $\mathcal{F}\to\min_\tau S(\tau)$（硬最小），$\lambda$ 大时趋于平均代价。
\end{definition}

\begin{note}
\textbf{配分函数难算，但算法从不显式算它.}\;$Z=\mathbb{E}_{\mathbb{Q}}[e^{-S/\lambda}]$ 是对整个轨迹空间的高维积分，看似无法计算（统计物理里算配分函数往往最难）。但 MPPI 最终要的不是 $Z$ 本身，而是吉布斯分布下的\emph{期望}（控制的加权均值），而期望是一个\emph{比值}：分子分母里的 $Z$ 完全相消。这正是后面“自归一化重要性采样”的精髓——权重 $w_k$ 的分母是用\emph{同一批样本}估计的 $Z$，在归一化里被约掉。自由能作为\emph{标量诊断量}很有用（它是要最小化的目标，也是风险敏感解读的对象），但算法\emph{从不显式计算 $Z$}。
\end{note}

\begin{theorem}[自由能--KL 对偶（Gibbs 变分原理）]\label{thm:mppi-duality}
自由能等于一个\emph{带 KL 正则的代价最小化}问题的最优值，
\begin{equation}
  \boxed{\ \mathcal{F}(\mathbb{Q})=\min_{\mathbb{P}}\ \Big\{\ \mathbb{E}_{\mathbb{P}}[S(\tau)]+\lambda\,\mathrm{KL}\big(\mathbb{P}\,\|\,\mathbb{Q}\big)\ \Big\}\ }
  \label{eq:mppi-dual}
\end{equation}
且最小值在\emph{吉布斯分布}
\begin{equation}
  \mathbb{P}^\star(\tau)=\frac{1}{Z}\,\mathbb{Q}(\tau)\,e^{-S(\tau)/\lambda},\qquad Z=\mathbb{E}_{\mathbb{Q}}[e^{-S/\lambda}]
  \label{eq:mppi-gibbs}
\end{equation}
处取到。
\end{theorem}
\begin{proof}
对任意与 $\mathbb{Q}$ 绝对连续的 $\mathbb{P}$，
\[
  \mathbb{E}_{\mathbb{P}}[S]+\lambda\,\mathrm{KL}(\mathbb{P}\|\mathbb{Q})
  =\lambda\,\mathbb{E}_{\mathbb{P}}\!\Big[\tfrac{S}{\lambda}+\log\tfrac{\mathrm{d}\mathbb{P}}{\mathrm{d}\mathbb{Q}}\Big]
  =\lambda\,\mathbb{E}_{\mathbb{P}}\!\Big[\log\tfrac{\mathrm{d}\mathbb{P}}{\mathrm{d}\mathbb{Q}\,e^{-S/\lambda}}\Big].
\]
把分母按 $\mathbb{Q}\,e^{-S/\lambda}=Z\,\mathbb{P}^\star$ 改写，得
\[
  =\lambda\,\mathbb{E}_{\mathbb{P}}\!\Big[\log\tfrac{\mathrm{d}\mathbb{P}}{\mathrm{d}\mathbb{P}^\star}\Big]-\lambda\log Z
  =\lambda\,\mathrm{KL}(\mathbb{P}\|\mathbb{P}^\star)-\lambda\log Z.
\]
因 $\mathrm{KL}(\mathbb{P}\|\mathbb{P}^\star)\ge0$，当且仅当 $\mathbb{P}=\mathbb{P}^\star$ 时为零，故右端在 $\mathbb{P}=\mathbb{P}^\star$ 取最小值 $-\lambda\log Z=\mathcal{F}(\mathbb{Q})$。
\end{proof}

\noindent 这个证明顺带暴露一个值得记住的事实：对偶式右端经改写恰好是 $\lambda\,\mathrm{KL}(\mathbb{P}\|\mathbb{P}^\star)-\lambda\log Z$。也就是说，\emph{最小化“期望代价 + KL 正则”等价于最小化 $\mathrm{KL}(\mathbb{P}\|\mathbb{P}^\star)$}——把候选分布 $\mathbb{P}$ 推向吉布斯目标 $\mathbb{P}^\star$。这里的 KL 写法是 $\mathrm{KL}(\mathbb{P}\|\mathbb{P}^\star)$（候选在前、目标在后），在变分推断里是\emph{反向 KL}，对应“模式寻找（mode-seeking）”的 I-投影。当 $\mathbb{P}$ 不受限时其最小值就是 $\mathbb{P}^\star$ 本身；一旦把 $\mathbb{P}$ 限制在可处理族里（如把控制更新限制成高斯），最小化反向 KL 就成了“用高斯去贴吉布斯目标主模式”的投影——这恰是 MPPI 与 CEM 分道扬镳之处（\cref{sec:mppi-family}）。

\begin{note}
\textbf{为什么偏偏是 KL？}\;一个自然的追问：正则项度量“$\mathbb{P}$ 离基分布 $\mathbb{Q}$ 多远”，为何用 KL 而非 $\chi^2$、Wasserstein？最实在的理由是\emph{它让问题有闭式解}——只有 $\lambda\mathrm{KL}(\mathbb{P}\|\mathbb{Q})$ 这种“对数密度比的期望”形式，配上期望代价，变分最优解才恰好是吉布斯分布 $\mathbb{Q}\,e^{-S/\lambda}$（指数族）；换成 $\chi^2$ 或 Wasserstein，最优解不再是干净的指数重加权，权重公式无从谈起。另一层理由：KL 正则对应信息论的最小相对熵原理、对应控制即推断的贝叶斯条件化（\cref{sec:mppi-inference}），还对应 \cref{sec:mppi-hjb} 那项二次控制代价的连续推广（Girsanov）。从 Kappen 到 Williams 到整个最大熵 RL，正则项清一色是 KL/相对熵——选它不是习惯，是数学逼着选的。
\end{note}

\paragraph{对偶式就是 SAC 的目标。}右端 $\mathbb{E}_{\mathbb{P}}[S]+\lambda\,\mathrm{KL}(\mathbb{P}\|\mathbb{Q})$ 是“期望代价 + 温度 $\times$ 对基分布的偏离惩罚”。把代价换成负奖励、把“对基分布的 KL”换成“对均匀分布的 KL（即负熵）”，这就\emph{逐字是 SAC 的最大熵目标}\cite{haarnoja2018soft}——奖励加上温度乘策略熵，最优解都是吉布斯/玻尔兹曼分布。这不是类比，是同一个变分原理在两个领域的两次应用（\cref{tab:mppi-sac}）。

\begin{table}[H]\centering\small
\caption{同一变分原理：最大熵 RL（SAC）与路径积分控制（MPPI）}
\label{tab:mppi-sac}
\begin{tabular}{lll}
\toprule
 & 最大熵 RL（SAC） & 路径积分控制（MPPI） \\
\midrule
决策对象 & 策略 $\pi(a\mid s)$ & 轨迹分布 $\mathbb{P}(\tau)$ \\
目标 & $\max\ \mathbb{E}[R]+\alpha\mathcal{H}(\pi)$ & $\min\ \mathbb{E}[S]+\lambda\,\mathrm{KL}(\mathbb{P}\|\mathbb{Q})$ \\
温度 & $\alpha$ & $\lambda$ \\
最优解 & 玻尔兹曼策略 $\pi^\star\propto e^{Q/\alpha}$ & 吉布斯分布 $\mathbb{P}^\star\propto\mathbb{Q}\,e^{-S/\lambda}$ \\
正则基准 & 均匀分布（熵） & 基分布 $\mathbb{Q}$（KL） \\
\bottomrule
\end{tabular}
\end{table}

\subsection{从吉布斯分布到 MPPI 权重公式}
\label{sec:mppi-weights}

我们拿到了最优轨迹分布 $\mathbb{P}^\star\propto\mathbb{Q}\,e^{-S/\lambda}$，但它是\emph{隐式}的——无法直接采样。MPPI 用重要性采样把它的均值近似出来。

\begin{theorem}[MPPI 权重公式与控制更新]\label{thm:mppi-weights}
从基分布 $\mathbb{Q}$ 采 $K$ 条 rollout（第 $k$ 条由噪声 $\boldsymbol{\varepsilon}_k$ 生成、总代价 $S_k$），任意 $\mathbb{P}^\star$ 下的期望可由\emph{自归一化重要性权重}估计：
\begin{equation}
  w_k=\frac{e^{-S_k/\lambda}}{\sum_{j=1}^{K}e^{-S_j/\lambda}},\qquad\sum_k w_k=1.
  \label{eq:mppi-w}
\end{equation}
把 $\mathbb{P}^\star$ 投影回参数化控制序列（$\mathbb{Q}$ 下第 $k$ 条轨迹时刻 $t$ 的控制是 $u_t+\varepsilon_k(t)$），$\mathbb{P}^\star$ 下的均值控制对应“把噪声按 $w_k$ 加权平均”，给出 MPPI 的核心更新
\begin{equation}
  \boxed{\ u_t\;\leftarrow\;u_t+\sum_{k=1}^{K}w_k\,\varepsilon_k(t)\ },\qquad t=0,\dots,T-1.
  \label{eq:mppi-update}
\end{equation}
\end{theorem}

\begin{theorem}[数值稳定化：$\rho$ 减法]\label{thm:mppi-rho}
softmax 对“所有指数同时减一个常数”不变。取 $\rho=\min_k S_k$，权重写作
\begin{equation}
  w_k=\frac{e^{-(S_k-\rho)/\lambda}}{\sum_j e^{-(S_j-\rho)/\lambda}},
  \label{eq:mppi-w-rho}
\end{equation}
其数值与 \cref{eq:mppi-w} 完全相同，但最小代价那条的指数变成 $e^0=1$、其余都在 $(0,1]$，既不上溢也不下溢。
\end{theorem}
\begin{proof}
给每个 $S_k$ 减同一常数 $\rho$，分子分母同乘 $e^{\rho/\lambda}$，归一化后约去：$\tfrac{e^{-(S_k-\rho)/\lambda}}{\sum_j e^{-(S_j-\rho)/\lambda}}=\tfrac{e^{\rho/\lambda}e^{-S_k/\lambda}}{e^{\rho/\lambda}\sum_j e^{-S_j/\lambda}}=w_k$。取 $\rho=\min_k S_k$ 使 $S_k-\rho\ge0$、$-(S_k-\rho)/\lambda\le0$、$e^{(\cdot)}\in(0,1]$，绝不上溢；最小代价那条恰得 $e^0=1$，是权重最大的一条。
\end{proof}

\noindent 直接算 $e^{-S_k/\lambda}$ 极易溢出：真实任务的 $S_k$ 动辄几百上千，取 $\lambda=1,S_k=1200$ 则 $e^{-1200}$ 在 $64$ 位浮点下硬下溢成 $0$，且\emph{所有}权重一起下溢、归一化 $0/0$ 给出 \texttt{NaN}，控制器当场失效。这不是小概率事件，而是“只要代价尺度大就必然发生”的确定性故障。\cref{eq:mppi-w-rho} 的 $\rho$ 减法即数值计算里历史悠久的 \textbf{log-sum-exp 技巧}——也是 \cref{sec:mppi-hjb} 那条本质洞察（$\min\leftrightarrow\log\sum\exp$）在算法层面的化身，是\emph{必需的}数值保命措施而非可选优化。\cref{sec:mppi-algorithm} 给出完整代码。

\subsection{权重公式的离散推导：高斯密度与控制代价为何消去}
\label{sec:mppi-weights-derivation}

\cref{thm:mppi-weights} 是从抽象的 $\mathbb{P}^\star\propto\mathbb{Q}\,e^{-S/\lambda}$ 写出的。两个常让人困惑的细节值得在离散情形显式算一遍：(1) 重要性权重里为什么没有高斯采样密度这一项？(2) \cref{sec:mppi-hjb} 明明有一项显式的二次控制代价 $\tfrac12 u^\top Ru$，到了权重里却不见了，它去哪了？

\begin{derivation}[高斯密度抵消与控制代价的再现]
\label{der:mppi-weights}
\textbf{(1) 高斯密度为何抵消.}\;把基分布写在扰动序列 $\boldsymbol{\varepsilon}=(\varepsilon_0,\dots,\varepsilon_{H-1})$ 上，每步 $\varepsilon_t\sim\mathcal{N}(\mathbf{0},\boldsymbol{\Sigma})$ 独立，密度 $q(\boldsymbol{\varepsilon})\propto\exp(-\tfrac12\sum_t\varepsilon_t^\top\boldsymbol{\Sigma}^{-1}\varepsilon_t)$。吉布斯最优分布 $\mathbb{P}^\star(\boldsymbol{\varepsilon})\propto q(\boldsymbol{\varepsilon})\,e^{-S(\boldsymbol{\varepsilon})/\lambda}$。要估计 $\mathbb{E}_{\mathbb{P}^\star}[\boldsymbol{\varepsilon}]$（用来更新 $U$），从 $q$ 采 $K$ 个样本做自归一化重要性采样，第 $k$ 个样本的权重
\[
  w_k\ \propto\ \frac{\mathbb{P}^\star(\boldsymbol{\varepsilon}_k)}{q(\boldsymbol{\varepsilon}_k)}\ \propto\ \frac{q(\boldsymbol{\varepsilon}_k)\,e^{-S_k/\lambda}}{q(\boldsymbol{\varepsilon}_k)}\ =\ e^{-S_k/\lambda}.
\]
$q$ 被精确约掉——因为我们恰好从那个出现在 $\mathbb{P}^\star$ 里的同一个 $q$ 采样。所以权重只剩代价的 softmax，没有任何高斯密度项。这正是 \cref{sec:mppi-algorithm} 代码只用代价、不碰采样密度就能算权重的原因。\\
\textbf{(2) 控制代价去哪了.}\;在最简 MPPI 里（围绕当前控制序列 $U$ 采\emph{零均值}扰动，即提议分布就是基分布 $\mathbb{Q}$），控制代价是\emph{隐式}的：它藏在“$\mathbb{P}^\star$ 不能离 $\mathbb{Q}$ 太远”这件事里（对偶式里的 $\lambda\mathrm{KL}(\mathbb{P}\|\mathbb{Q})$ 项），你不必手动往 $S$ 里加 $\tfrac12\varepsilon^\top R\varepsilon$。但一旦你想\emph{移动提议分布}（围绕一个不同于 $U$ 的均值、或用不同协方差采样——Williams 2017 的 generalized importance sampling\cite{williams2017model}），权重就要乘上提议分布 $q'$ 与基分布 $q$ 的密度比 $w_k\propto e^{-S_k/\lambda}\cdot q(\boldsymbol{\varepsilon}_k)/q'(\boldsymbol{\varepsilon}_k)$，其中 $\log(q/q')$ 是控制扰动的二次型。两个高斯之间的对数密度比恰是一个 $\propto u^\top\boldsymbol{\Sigma}^{-1}u$ 的二次型——而在 Kappen 关系 $\boldsymbol{\Sigma}=\lambda R^{-1}$ 下，$u^\top\boldsymbol{\Sigma}^{-1}u=u^\top Ru/\lambda$，于是这个换测度修正在权重里正是 $e^{-(\frac12 u^\top Ru)/\lambda}$——\emph{恰好就是} \cref{sec:mppi-hjb} 那项二次控制代价进入指数后的样子。
\end{derivation}

\noindent 换句话说：\cref{sec:mppi-hjb} 显式写出的二次控制代价，与 \cref{sec:mppi-freeenergy} 对偶式里的 KL 项，是\emph{同一个东西}——受控路径测度与被动路径测度之间的对数密度比（连续时间里这叫 Girsanov 变换）。最简 MPPI 把提议分布选成基分布，于是这一项退化为常数、权重纯由代价决定；只有当你移动提议分布时，它才作为显式的二次项重新现身——这就把 \cref{sec:mppi-hjb} 的 PDE 推导与 \cref{sec:mppi-freeenergy} 的信息论推导在这一点上彻底缝合。

\subsection{不可微代价的兑现}
\label{sec:mppi-nondiff}

\cref{sec:mppi-freeenergy} 反复强调“$S_k$ 只作一个标量进指数”，所以代价可以不连续、不可微、来自黑箱。这一点值得亲眼一见，因为它正是 MPPI 相对梯度法最硬的优势。设一个\emph{阶梯型}终端代价：终点落入目标带 $[1.5,2.5]$ 代价为 $0$，否则代价 $10$——这个代价的梯度\emph{处处为零}，梯度下降在它上面完全收不到信号、原地不动。但 MPPI 的权重不在乎可微性：

\begin{codebox}[Python]{不可微（阶梯）代价:指数权重照样工作}
import numpy as np
rng = np.random.default_rng(1)
dt, H, sigma = 0.05, 40, 1.5; K = 8000
eps = sigma*rng.standard_normal((K,H)); U = eps          # x0=0, nominal control zero
x = np.zeros(K); S = np.zeros(K)
for t in range(H):
    x += U[:,t]*dt; S += 0.5*0.05*U[:,t]**2*dt
S += np.where((x>=1.5)&(x<=2.5), 0.0, 10.0)              # step terminal cost (non-smooth)
w = np.exp(-(S - S.min())); w /= w.sum()
print(f"in-band frac={((x>=1.5)&(x<=2.5)).mean()*100:.1f}%  in-band weight={w[(x>=1.5)&(x<=2.5)].sum():.3f}")
print(f"weighted x_T={(w*x).sum():+.3f}   raw mean x_T={x.mean():+.3f}")
\end{codebox}

\noindent 输出为“落入目标带占比 $=0.1\%$、带内总权重 $=0.955$”与“加权 $x_T=+1.53$、裸平均 $x_T=-0.009$”：尽管只有 $0.1\%$ 的 rollout 侥幸落进目标带，权重却把 $95.5\%$ 的信任压在它们身上，加权终点 $+1.53$ 稳稳落在带内；不加权的裸平均还停在出发点附近。\textbf{指数权重 $e^{-S/\lambda}$ 只需比较代价的大小，从不对代价求导}，所以阶梯、障碍的硬跳变、黑箱仿真器吐出的不可微代价，对它都一视同仁。这正是 MPPI 能直接吃碰撞惩罚、稀疏奖励、规则化约束这些让梯度法束手无策的代价的原因。

\subsection{MPPI 与 REINFORCE 的同构}
\label{sec:mppi-reinforce}

把 MPPI 更新和你熟悉的 REINFORCE 并排看。REINFORCE 对高斯策略 $\pi=\mathcal{N}(u,\boldsymbol{\Sigma})$ 的策略梯度上升是（$R=-S$ 为负代价，$b$ 为基线）\cite{williams1992simple}：
\begin{equation}
  u\;\leftarrow\;u+\eta\,\mathbb{E}\big[(R-b)\,\nabla_u\log\pi\big]=u+\eta\,\mathbb{E}\big[-(S-b)\,\boldsymbol{\Sigma}^{-1}\varepsilon\big],
  \label{eq:mppi-reinforce}
\end{equation}
而 MPPI 更新是 $u\leftarrow u+\mathbb{E}_{\mathbb{Q}}[\,(e^{-S/\lambda}/Z)\cdot\varepsilon\,]$。\textbf{两者骨架完全相同}：都是 $u\leftarrow u+\mathbb{E}[\,g(S)\cdot\varepsilon\,]$——“把采样噪声 $\varepsilon$ 按一个与代价相关的标量权重 $g(S)$ 加权平均”。区别只在权重函数 $g$（\cref{tab:mppi-reinforce}）。

\begin{table}[H]\centering\small
\caption{MPPI 与 REINFORCE：同一骨架的两种权重函数}
\label{tab:mppi-reinforce}
\begin{tabular}{llll}
\toprule
 & 权重函数 $g(S)$ & 性质 & 需要的东西 \\
\midrule
REINFORCE & $\propto-(S-b)$（\textbf{线性}） & 一阶；可正可负 & 学习率 $\eta$、基线 $b$、$\boldsymbol{\Sigma}^{-1}$ \\
MPPI & $\propto e^{-S/\lambda}$（\textbf{指数 / softmax}） & 零阶；恒非负、自归一化 & 温度 $\lambda$；无学习率、无矩阵逆 \\
\bottomrule
\end{tabular}
\end{table}

\noindent 二者的联系可说得更精确：把 MPPI 的指数权重在 $1/\lambda$ 小时做一阶泰勒展开 $e^{-S/\lambda}\approx1-S/\lambda$，归一化后权重 $\propto-(S-\bar S)$——\emph{正是带基线的 REINFORCE}。

\begin{insight}[MPPI = 零阶策略梯度的闭式特例；为什么没有学习率]
MPPI 是\textbf{零阶（zeroth-order）策略梯度的闭式特例}，或说“softmax 版的 REINFORCE”。指数权重 $e^{-S/\lambda}$ 是 KL 正则最优化问题（\cref{thm:mppi-duality}）的\emph{精确闭式解}；REINFORCE 的线性权重则是它的\emph{一阶泰勒近似}。更深一层：MPPI 的更新是分布空间里的一步\emph{自然梯度 / 镜像下降}，而步长被 KL 约束自动定死。普通梯度下降在\emph{参数}的欧氏空间里走，学习率是要手调的自由量；但当优化对象是一个\emph{概率分布}时，两个分布“差多远”该用 KL 度量，对应 Fisher 信息度量下的\emph{信息几何}，在其中做最速下降就是自然梯度。\cref{thm:mppi-duality} 的对偶问题恰是一个\emph{带 KL 信赖域的最优化}：“在离当前分布 $\mathbb{Q}$ 不太远（KL 惩罚）的前提下把期望代价降到最低”，它有闭式解（吉布斯分布），于是 MPPI 不需沿梯度试探着走——它直接跳到信赖域内的最优分布。所谓“没有学习率”，其实是\emph{学习率被 KL 信赖域 + 闭式解吸收了}。$\lambda$ 同时扮演信赖域半径：$\lambda$ 大、信赖域小、每步挪得保守；$\lambda$ 小、信赖域大、挪得激进。这把本章的几条线和更广的优化版图缝在一起——CMA-ES 的现代推导、自然进化策略（NES）、MPO 的 E 步，走的都是“在 KL/信息几何信赖域里更新分布”这同一条路。
\end{insight}

\paragraph{收敛性的边界。}把 MPPI 的迭代看成镜像下降：每一步在“离当前分布 KL 不太远”的信赖域里跳到期望代价更低的吉布斯分布。在\emph{理想（无限样本）}情形，这一步是 majorize-minimize / EM 式更新，\emph{单调不增}地降低 KL 正则后的目标——这就是实测里代价稳步下降的理论来源，和 MPO 的 E 步、CEM 的迭代收敛性是同一类论证。但有两条边界：\emph{其一，有限样本带来随机性}——真实权重是 $K$ 条 rollout 的蒙特卡洛估计（有偏、有方差，见 \cref{sec:mppi-family}），单调下降只在期望/大 $K$ 下成立，小 $K$ 时单步可能不降反升，迭代是一个\emph{随机}逼近过程；\emph{其二，非凸代价只保证局部}——信赖域里的下降把分布推向\emph{当前邻域}内的更优解，代价多模时收敛到\emph{某一个}模式而非全局最优。诚实的表述是：MPPI 在理想条件下单调改善 KL 正则目标、收敛到一个局部最优分布；有限样本让它变成随机逼近，非凸让“最优”退化为“局部最优”。这在工程上完全够用——配合高频重规划，每周期只需把控制朝更优挪一步，并不指望单次迭代求到全局最优。这条更紧的、二次代价下的收敛速率结论由 CoVO-MPC 给出，见 \cref{sec:mppi-family}。

\begin{pitfall}[思维陷阱：把 MPPI 更新当成“梯度下降走一步”]
因为它和 REINFORCE 同构，就以为 MPPI 也是“沿某个梯度方向、按学习率走一步”。\textbf{后果}：去找 MPPI 的“学习率”并试图调它，或担心“步长太大不收敛”，套用梯度法的收敛分析。\textbf{根本原因}：MPPI 解的是\emph{闭式}最优分布（吉布斯分布）的投影——它是\emph{加权平均}，不是梯度步；没有学习率、没有矩阵逆（这正是 PI$^2$ 的标志性卖点）。\textbf{对策}：把每次更新理解为“用当前这批样本，对最优分布的均值做一次蒙特卡洛估计”。控制“激进程度”的旋钮是温度 $\lambda$（权重分布锐度）与采样协方差 $\boldsymbol{\Sigma}$（探索几何范围），不是学习率。
\end{pitfall}

\begin{pitfall}[概念误区：以为采样数 $K$ 越大就一定越好]
性能不好就一味加大 $K$。\textbf{后果}：$K$ 加到很大，显存/算力吃满，效果却仍差——问题根本不在样本数。\textbf{根本原因}：$K$ 大只降低\emph{估计方差}；但若基分布 $\mathbb{Q}$ 离最优分布 $\mathbb{P}^\star$ 太远（初始控制序列很糟、协方差不合适），低代价区域\emph{根本没有样本落进去}，再多样本也是在错误的地方平均（\cref{sec:mppi-mc} 的“提议分布不重叠”问题）。\textbf{对策}：先保证采样分布质量——合理的 warm-start（\cref{sec:mppi-algorithm}）、合适的协方差 $\boldsymbol{\Sigma}$、必要时注入先验。\emph{分布质量优先于样本数}。
\end{pitfall}

\begin{practice}
\begin{enumerate}
  \item 用拉格朗日乘子法独立证明对偶定理：在约束 $\int\mathbb{P}(\tau)\,\mathrm{d}\tau=1$ 下最小化 $\mathbb{E}_{\mathbb{P}}[S]+\lambda\mathrm{KL}(\mathbb{P}\|\mathbb{Q})$，导出最优解为吉布斯分布 \cref{eq:mppi-gibbs}，并代回验证最优值恰为 $-\lambda\log Z=\mathcal{F}(\mathbb{Q})$。与正文的 KL 配方证法对照。
  \item 从 \cref{eq:mppi-gibbs} 出发，写出“用 $K$ 个 $\mathbb{Q}$ 样本估计 $\mathbb{P}^\star$ 下控制均值”的自归一化重要性采样估计量，导出权重 \cref{eq:mppi-w} 与更新 \cref{eq:mppi-update}。说明为何“自归一化”（分母用样本和而非真实 $Z$）让我们\emph{不需要知道配分函数 $Z$}。
  \item 证明自由能的均值-方差展开 $\mathcal{F}=\mathbb{E}[S]-\tfrac{1}{2\lambda}\mathrm{Var}[S]+O(1/\lambda^2)$（把指数与对数都对 $1/\lambda$ 泰勒展开），并据此说明 $\lambda\to\infty$ 时风险中性、$\lambda$ 有限时偏好高方差（见 \cref{sec:mppi-temperature} 风险敏感）。
  \item 证明自由能的双侧夹逼 $\min_\tau S(\tau)\le\mathcal{F}(\mathbb{Q})\le\mathbb{E}_{\mathbb{Q}}[S]$（左用 $\mathbb{E}[e^{-S/\lambda}]\le e^{-\min S/\lambda}$，右用 Jensen 于凹函数 $\log$）。再证次优性证书恒等式：对任意候选 $\mathbb{P}$，$\mathbb{E}_{\mathbb{P}}[S]+\lambda\mathrm{KL}(\mathbb{P}\|\mathbb{Q})=\mathcal{F}(\mathbb{Q})+\lambda\mathrm{KL}(\mathbb{P}\|\mathbb{P}^\star)$。说说这两条分别怎么用作调试 MPPI 的健全性检查与“离最优多远”的度量。
\end{enumerate}
\end{practice}

% ════════════════════════════════════════════════════════════════
\section{控制即推断：MPPI、SAC、REINFORCE 的共同母体}
\label{sec:mppi-inference}

读到这里，一个问题应已浮现：为什么 MPPI 的权重、SAC 的玻尔兹曼策略、LMDP 的 desirability、乃至 REINFORCE 的更新，\emph{全都长着 $e^{-\text{代价}/\lambda}$ 这副面孔}？这不是巧合——它们是同一个框架\textbf{控制即推断（control as inference）}的不同侧影。这个视角由 Levine 2018 的综述系统化\cite{levine2018reinforcement}；更早地，Kappen、Gómez、Opper 2012 已证明 Todorov 的 KL 控制理论“把路径积分控制作为一个特例包含其中”\cite{kappen2012optimal}。

\begin{definition}[最优性变量与规划即推断]\label{def:mppi-inference}
在每个时刻引入一个二值\emph{最优性变量}（optimality variable）$\mathcal{O}_t\in\{0,1\}$，规定它取 $1$ 的概率随该步代价指数衰减：$p(\mathcal{O}_t{=}1\mid x_t,u_t)\propto\exp(-c(x_t,u_t)/\lambda)$，整条轨迹代价 $S(\tau)=\sum_t c(x_t,u_t)$（含终端）。假设“观测到”整条轨迹自始至终都最优（$\mathcal{O}_{1:T}{=}1$），轨迹的后验由贝叶斯公式给出
\begin{equation}
  p(\tau\mid\mathcal{O}_{1:T}{=}1)\ \propto\ \underbrace{p(\tau)}_{\mathbb{Q}(\tau)}\;\prod_t p(\mathcal{O}_t{=}1\mid x_t,u_t)\ \propto\ \mathbb{Q}(\tau)\,e^{-S(\tau)/\lambda}.
  \label{eq:mppi-posterior}
\end{equation}
\end{definition}

\noindent \textbf{这个后验，恰恰就是 \cref{thm:mppi-duality} 的吉布斯分布 $\mathbb{P}^\star$。}也就是说——“最优轨迹分布”与“在观测到全程最优这一证据后、对轨迹做贝叶斯推断得到的后验”是同一个东西。\textbf{规划（planning）就是推断（inference）。}那个反复出现的自由能立刻有了概率意义：$\mathcal{F}(\mathbb{Q})=-\lambda\log\mathbb{E}_{\mathbb{Q}}[e^{-S/\lambda}]=-\lambda\log p(\mathcal{O}_{1:T}{=}1)$，即自由能等于“全程最优”这一证据的对数边际似然乘以 $-\lambda$。于是对偶定理 $\mathcal{F}=\min_{\mathbb{P}}\{\mathbb{E}_{\mathbb{P}}[S]+\lambda\mathrm{KL}(\mathbb{P}\|\mathbb{Q})\}$ 正是\emph{变分推断}的标准形式：右端是负的证据下界（ELBO），$\lambda\mathrm{KL}(\mathbb{P}\|\mathbb{Q})$ 是变分分布与真后验的间隙，$\mathbb{P}=\mathbb{P}^\star$（真后验）时间隙归零、下界紧贴对数证据。\cref{sec:mppi-freeenergy} 整套推导，本质上就是一次变分推断。

\begin{table}[H]\centering\small
\caption{四条路线，同一个后验 $\mathbb{P}^\star\propto\mathbb{Q}\,e^{-S/\lambda}$}
\label{tab:mppi-inference}
\begin{tabularx}{\linewidth}{lLl}
\toprule
路线 & 怎么算后验 $\mathbb{P}^\star\propto\mathbb{Q}\,e^{-S/\lambda}$ & 出处 \\
\midrule
Kappen 路径积分 & 连续时间，指数变换 + Feynman--Kac，对被动轨迹采样 & \cref{sec:mppi-hjb,sec:mppi-fk} \\
Todorov LMDP & 离散时间，值函数路线，线性贝尔曼 $z=MPz$ & \cref{sec:mppi-lmdp} \\
MPPI & 轨迹空间，重要性采样\emph{在线}估计后验均值；每步重规划、不学策略 & \cref{sec:mppi-freeenergy} \\
SAC / soft Q / MPO & 同一后验，但把推断\emph{摊销（amortize）}进一张学到的策略网络；离线训练、可复用 & \cite{haarnoja2018soft,abdolmaleki2018maximum} \\
\bottomrule
\end{tabularx}
\end{table}

\noindent 最后一行点出 MPPI 与 SAC 最干净的区别：\emph{两者解的是同一个 KL 正则推断问题}；SAC 把答案摊销成一个离线学好、随用随取的策略，MPPI 则在每个控制周期\emph{在线}重算一遍（不学策略、不复用）。这也回答了“有了 SAC 为什么还要 MPPI”——当你手里有一个（哪怕黑箱的）模型、又负担得起在线算力时，MPPI 免去训练、即插即用、且能随时换代价函数；当你有海量离线交互、想要零在线开销的反应式策略时，SAC 的摊销更划算。不止 SAC——MPO（Maximum a Posteriori Policy Optimisation）把策略优化写成 EM，E 步算一个非参数的玻尔兹曼后验 $q\propto e^{Q/\eta}$（与 MPPI 权重同形）、M 步把它拟合进策略网络\cite{abdolmaleki2018maximum}；AWR（Advantage-Weighted Regression）用 $e^{A/\beta}$（优势的指数）给经验动作加权回归\cite{peng2019advantage}。四者的权重都是“某种价值量除以温度再取指数”，都来自同一个 KL 正则目标的闭式解，差别只在\emph{价值量是什么}（轨迹代价 / 优势 / Q 值）与\emph{怎么消费权重}（在线挪序列 / 回归 / 蒸馏成网络）。

\begin{insight}[一个母体，按资源约束做出的工程选择]
\cref{sec:mppi-reinforce} 那两个同构（MPPI $\approx$ REINFORCE、MPPI $\leftrightarrow$ SAC）之所以成立，根子全在这里：\emph{它们都是控制即推断的实例}。指数化代价后验 $\mathbb{P}^\star$ 是共同的“母体”，具体算法只是在回答同一问题“如何逼近这个后验”时，选择了不同的\emph{摊销程度}与\emph{计算路线}：有模型 + 在线算力 $\to$ MPPI（每步重算）；要可复用的反应式策略 + 离线交互 $\to$ SAC（摊销成网络）；小规模离散模型 $\to$ LMDP（线性代数一次解出）。看穿这一点，你就不会把 MPPI、SAC、LMDP 当成三门要分别死磕的手艺，而是同一框架下按资源约束做出的三种工程选择——你在任意一边积累的直觉，都能照亮另外两边。
\end{insight}

% ════════════════════════════════════════════════════════════════
\section{温度 \texorpdfstring{$\lambda$}{lambda}、风险敏感与对称性破缺}
\label{sec:mppi-temperature}

\paragraph{动机：唯一的自由参数，初学者几乎总答错。}
权重公式 $w_k\propto e^{-S_k/\lambda}$ 里，$\lambda$ 是除采样协方差外唯一的旋钮。但“该取多大”这个问题，初学者要么沿用某个默认值（“就用 $1.0$ 吧”），要么越调越小以为能逼出“更优解”——两种做法都会出事。要正确用 MPPI，必须先建立一个准确直觉：\textbf{$\lambda$ 控制的是“拿到一批样本后，你对它们有多挑剔”}。

\paragraph{反面：两个极端都退化成没用的东西。}$\lambda\to0$ 时 $w_k$ 坍缩成 $\delta$ 函数，只信代价最低的那\emph{一条} rollout——这条很可能只是某个噪声样本的侥幸，控制因此剧烈抖动、被单条坏样本带偏。$\lambda\to\infty$ 时 $w_k\to1/K$，对所有轨迹一视同仁地平均——好坏不分，相当于完全没用上代价信息，控制趋于“原地不动的平均”。两个极端之间那段“恰到好处”才是 MPPI 真正工作的区间。

\begin{theorem}[$\lambda$ 的三种极限与软硬最小值]\label{thm:mppi-lambda}
自由能的样本估计 $-\lambda\log\frac1K\sum_k e^{-S_k/\lambda}$ 是一个由 $\lambda$ 调节锐度的\emph{软最小值}（soft-min）。两端极限把 $\lambda$ 的作用一览无余（\cref{tab:mppi-lambda}）：$\lambda\to0$ 时指数把最小项之外全部压成零（硬 $\arg\min$）；$\lambda\to\infty$ 时对 $e^{-S/\lambda}\approx1-S/\lambda$ 做一阶展开，软最小值退化为算术平均。
\end{theorem}

\begin{table}[H]\centering\small
\caption{温度 $\lambda$ 的三种极限}
\label{tab:mppi-lambda}
\begin{tabular}{lllll}
\toprule
$\lambda$ & 权重行为 & 软最小值极限 & 策略含义 & SAC 对应 \\
\midrule
$\lambda\to0$ & $w_k\to\delta(k{=}\arg\min_j S_j)$ & $\to\min_k S_k$（硬最小） & 只信代价最低一条，贪心 & $\alpha\to0$（纯利用） \\
适中 & 权重平滑集中在低代价轨迹上 & 介于二者之间 & 一簇好轨迹上加权融合 & 最优 $\alpha$ \\
$\lambda\to\infty$ & $w_k\to1/K$ & $\to\frac1K\sum_k S_k$（平均） & 好坏不分，盲目平均 & $\alpha\to\infty$（纯探索） \\
\bottomrule
\end{tabular}
\end{table}

\paragraph{定量选取：用有效样本数（ESS）当探针。}与其凭感觉调 $\lambda$，不如盯住一个可观测的诊断量——\emph{有效样本数}
\begin{equation}
  \mathrm{ESS}=\frac{1}{\sum_k w_k^2}\ \in[1,K],
  \label{eq:mppi-ess}
\end{equation}
它衡量“名义上撒了 $K$ 条，实际有多少条在起作用”。$\lambda$ 太小、权重几乎全压在一两条上时 $\mathrm{ESS}\to1$（名义撒 $K$、实际只一条在工作，估计自然不准）；$\lambda$ 太大、权重趋于均等时 $\mathrm{ESS}\to K$（好坏不分）。健康区间通常取 $\mathrm{ESS}$ 为 $K$ 的若干分之一。这与 \cref{sec:mppi-mc} 的方差爆炸是同一件事的两种说法——纯被动采样时落在低代价区的样本极少、$\mathrm{ESS}$ 塌到接近 $1$，而好的提议分布把权重摊开、把 $\mathrm{ESS}$ 拉回健康比例。可进一步做 SAC 式的\emph{按目标自动调温}：设一个目标 ESS，在线根据当前 ESS 偏离方向微调 $\lambda$（ESS 太低升温、太高降温），把“ESS 仪表盘”从只读指标升级成闭环调节器——与 SAC 按目标熵调 $\alpha$ 是同一套思路。

\subsection{风险敏感视角：自由能其实是“指数效用”}
\label{sec:mppi-risk}

控制即推断把自由能解读成对数证据；还有一个对工程更直接的解读——\textbf{自由能就是控制论里的风险敏感目标}，它把 $\lambda$ 从“温度/锐度旋钮”进一步揭示为一个\emph{风险态度旋钮}。

\begin{theorem}[自由能的均值-方差展开]\label{thm:mppi-meanvar}
把自由能在大 $\lambda$ 处对 $1/\lambda$ 泰勒展开，得
\begin{equation}
  \mathcal{F}(\mathbb{Q})=-\lambda\log\mathbb{E}_{\mathbb{Q}}\big[e^{-S/\lambda}\big]\ \approx\ \mathbb{E}_{\mathbb{Q}}[S]\ -\ \frac{1}{2\lambda}\,\mathrm{Var}_{\mathbb{Q}}[S]\ +\ \cdots
  \label{eq:mppi-meanvar}
\end{equation}
\end{theorem}

\noindent 自由能不是平均代价，而是\emph{平均代价减去一个方差项}。这立刻给出 $\lambda$ 的风险含义：因为方差项前是减号，最小化 $\mathcal{F}$ 等于在压低平均代价的同时\emph{奖励代价的高方差}——偏爱那些“存在极低代价尾巴”的分布。换句话说，标准路径积分 / MPPI（$\lambda>0$）天生是\textbf{轻度冒险（risk-seeking）}的：指数权重把赌注押在最好的那几条 rollout 上。$\lambda\to0$ 是极端冒险（只认单条最优）；$\lambda\to\infty$ 时方差项消失、$\mathcal{F}\to\mathbb{E}[S]$，退回\emph{风险中性}的期望代价。这个目标在控制论里有专名：形如 $\tfrac{1}{\theta}\log\mathbb{E}[e^{\theta X}]$ 的量叫\emph{熵风险测度}（entropic risk measure），等价于\emph{指数效用下的确定性等价}，最早由 Jacobson 1973 引入随机控制（线性-指数-二次-高斯，LEQG）\cite{jacobson1973optimal}。van den Broek、Wiegerinck、Kappen 2010 证明：在与本章完全相同的假设下，路径积分直接推广到风险敏感控制——最小化的正是指数加权 cost-to-go，指数权重的符号与大小决定 risk-seeking 还是 risk-averse 的行为\cite{vandenbroek2010risk}。

\begin{insight}[$\lambda$ 的三张面孔是同一个量]
$\lambda$ 有三张面孔，对应本章三个视角，但它们是同一个量。在 \cref{sec:mppi-hjb,sec:mppi-fk} 它是\emph{噪声温度}（Kappen 条件 $\boldsymbol{\Sigma}=\lambda R^{-1}$ 把它绑定到控制代价）；在 \cref{sec:mppi-temperature} 它是 softmax 的\emph{锐度}（$\min$ 与平均之间的旋钮）；在这里它是\emph{风险态度}（冒险与中性之间的旋钮）。三者不是巧合，而是同一个指数权重 $e^{-S/\lambda}$ 在不同问题里的不同投影：温度调节探索的随机性、锐度调节权重的集中度、风险调节对代价分布尾部的偏好——而它们由同一个 $\lambda$ 一并决定。所以调 $\lambda$ 从来不是调一个孤立的超参数，而是同时在拨动探索、集中度与风险三个耦合的旋钮；理解这一点，“$\lambda$ 调不好”就从“性能差一点”升级为“行为性质、安全裕度一起变”的判断。
\end{insight}

\begin{pitfall}[思维陷阱：在随机环境里把 $\lambda$ 调小只是为了“更精准”]
认为温度越低控制就越接近“真正的最优”，于是在任何环境都倾向把 $\lambda$ 往小调。\textbf{后果}：在噪声大、模型不准的真实系统上，低温控制器表现\emph{激进而脆弱}——它押注于“一切顺利时能达到的最低代价”，遇扰动或建模误差就大幅偏离，甚至失稳。\textbf{根本原因}：自由能 $\approx\mathbb{E}[S]-\mathrm{Var}[S]/2\lambda$，$\lambda$ 同时是\emph{风险旋钮}；$\lambda$ 越小越偏好高方差（冒险），把信任压在代价分布最乐观的尾巴上。这在确定性环境无害，在随机环境就是鲁莽。\textbf{对策}：在随机/不确定环境里把 $\lambda$ 理解为“风险态度”而非单纯“锐度”——需要保守时不要一味降温；若任务本质要求 risk-averse，应转向风险敏感推广里 risk-averse 的那一侧（指数权重反号），而非默认的冒险方向。
\end{pitfall}

\begin{pitfall}[概念误区：混淆温度 $\lambda$ 与采样协方差 $\boldsymbol{\Sigma}$ 的作用]
以为 $\lambda$ 与 $\boldsymbol{\Sigma}$ 都是“控制探索”的、调哪个都一样。\textbf{后果}：想让 MPPI 更激进却拧错旋钮，或把两者耦合调坏。\textbf{根本原因}：$\boldsymbol{\Sigma}$ 决定\emph{采样在控制空间撒得多开}（探索的几何范围）；$\lambda$ 决定\emph{拿到样本后对它们多挑剔}（权重分布的锐度）。一个管“撒哪里”，一个管“信哪条”。\textbf{对策}：先用 $\boldsymbol{\Sigma}$ 确保采样覆盖到有意义的控制范围，再用 $\lambda$ 调节“在这些样本里多贪心地偏向低代价者”。（\cref{sec:mppi-family} 的 CoVO-MPC 进一步指出二者在收敛率里其实耦合，求最优 $\boldsymbol{\Sigma}$ 是在给定 $\lambda$ 下求的。）
\end{pitfall}

\begin{practice}
\begin{enumerate}
  \item 严格证明 \cref{thm:mppi-lambda} 的两个极限：$\lambda\to0$ 时软最小值 $\to\min_k S_k$（指数把非最小项压零）；$\lambda\to\infty$ 时 $\to\frac1K\sum_k S_k$（一阶展开）。
  \item 推导 \cref{eq:mppi-meanvar}：把 $\mathcal{F}=-\lambda\log\mathbb{E}[e^{-S/\lambda}]$ 的指数与对数都对 $1/\lambda$ 泰勒展开到二阶。再思考：若把目标改成 $+\lambda\log\mathbb{E}[e^{+S/\lambda}]$（指数权重反号），方差项变号，得 risk-averse 还是 risk-seeking？（这正是 van den Broek 等 2010 调节风险态度的机制。）
  \item 实现一个温度自适应：设目标 ESS，在线根据当前 $\mathrm{ESS}=1/\sum_k w_k^2$ 偏离目标的方向微调 $\lambda$，在一个到达任务上观察 $\lambda$ 是否自动收敛到一个合理量级。
\end{enumerate}
\end{practice}

% ════════════════════════════════════════════════════════════════
\section{接收水平 MPPI 算法}
\label{sec:mppi-algorithm}

\cref{thm:mppi-weights,thm:mppi-rho} 给出了 MPPI 的数学核心：采样、按 softmax 加权、更新控制序列。但要把它落成能在真实机器人上以毫秒级运行的控制器，还需把单次开环求解嵌进\emph{接收水平（receding horizon）}的闭环，并补两个工程细节。本节把理论核补成可部署的算法骨架（GPU 并行、Savitzky--Golay 平滑等纯工程细节属感知/控制工程范畴，本理论章只点到为止）。

\begin{algo}[接收水平 MPPI（可部署骨架）]\label{alg:mppi}
单次求解给出的是一条\emph{开环}控制序列；“执行第一步—重测—重规划”的循环把它变成\emph{闭环反馈}。这正是所有 MPC 的共性，也是与 \cref{sec:plan-traj} 离线轨迹优化的根本区别。
\end{algo}

\begin{codebox}[text]{Information-Theoretic MPPI（可部署版伪代码）}
INPUT:  init state x0, current control sequence U={u_0..u_{T-1}},
        temperature lambda, sample count K, covariance Sigma
OUTPUT: executed control u_0, shifted sequence U' (warm-start for next cycle)

each control cycle:
  // -- 1. parallel sampling and rollout (independent across k => GPU) --
  for k = 1..K  in parallel:
    eps_k ~ N(0, Sigma)              // noise sequence of length T
    x <- x0;  S_k <- 0
    for t = 0..T-1:
      v_t = u_t + eps_k(t)           // perturbed control
      x <- F(x, v_t)                 // forward dynamics (black box: sim / NN)
      S_k += running_cost(x, v_t)    // accumulate running cost  (NOT optional!)
    S_k += terminal_cost(x)
  // -- 2. weights (numerically stable, Theorem rho-subtraction) --
  rho = min_k S_k                    // log-sum-exp baseline
  for k = 1..K:  w_k = exp(-(S_k - rho)/lambda)
  eta = sum_k w_k;  w_k <- w_k / eta
  // -- 3. weighted update of control sequence --
  for t = 0..T-1:  u_t <- u_t + sum_k w_k * eps_k(t)
  // -- 4. (optional) Savitzky-Golay smoothing to tame chattering --
  U <- smooth(U)
  // -- 5. execute u_0, warm-start shift (reuse this cycle's work) --
  execute(u_0)
  U' <- {u_1, u_2, ..., u_{T-1}, u_init}     // shift left, pad tail
\end{codebox}

\paragraph{接收水平：单次开环 $\to$ 闭环反馈。}MPPI 每周期算出一整条长度 $T$ 的控制序列，但只执行第一个 $u_0$，下一周期在\emph{新测得}的状态上重新规划。为什么不一次算好整条就盲目执行？两个原因，与所有 MPC 共通：其一，\emph{模型不完美}——用模型 rollout 算出的“最优序列”基于模型预测，真实系统执行后会偏离（扰动、模型失配），一次算好就盲执整条会让偏差累积到不可收拾；接收水平的精髓是“只执行第一步、用真实新状态重规划”，不断用最新反馈纠正预测偏差。其二，\emph{环境在变}——障碍在动、目标在变，一次规划无法预见未来所有变化，必须持续重规划。时域 $T$ 是一个权衡：太短则远见不足（撞上视野外的障碍），太长则采样维度 $T\times m$ 升高、好序列更难撞中（\cref{sec:mppi-mc} 的高维采样困难）。

\begin{insight}[warm-start 就是在构造一个好的提议分布]
执行 $u_0$、系统前进一步后，\emph{剩下的 $(u_1,\dots,u_{T-1})$ 仍然几乎最优}——机器人只动了一步，“从下一状态出发的最优计划”与“上周期算出的、去掉第一步的计划”高度重合。warm-start 把整条序列\emph{左移一位}、末位补一个默认控制 $u_{\text{init}}$，作为下周期初值。这不只是“省点计算”：\cref{sec:mppi-freeenergy} 证明过，重要性采样的理想（零方差）提议分布就是目标分布本身，MPPI 的全部改进本质都是“让提议分布更接近那个隐式最优分布”。warm-start 正是这一原则最廉价、最有效的实现——上周期优化出的控制序列已是对“最优控制长什么样”的相当好的估计，以它为采样中心，等于把提议分布直接挪到最优解附近。这解释了一个乍看反直觉的现象：warm-start 的威力常\emph{大于}单纯把 $K$ 翻倍——前者改善提议分布的“位置”（采样撒在哪），后者只增加“样本数”（原地多撒几个），而位置比数量重要得多。
\end{insight}

\paragraph{两个工程细节，不做就不能用。}其一，$\rho$ 减法（\cref{thm:mppi-rho}）——伪代码第 2 步，不做则代价尺度一大就全体下溢成 \texttt{NaN}、控制器失效。其二，warm-start shift——第 5 步，不做则每周期从零优化、高频控制下根本算不完。还有一个执行器层面的细节：原始 MPPI 的输出是大量随机轨迹的加权平均，天然带高频噪声，直接喂电机会嗡鸣磨损，故第 4 步常加 Savitzky--Golay 平滑\cite{savitzky1964smoothing}（信号处理里的多项式滑窗平滑，本是平滑光谱数据的）。末位填充也有讲究：填零意味着“假设末端不施加控制”（对“减速停止”类任务合理），但对需持续控制才能维持的系统（如悬停四旋翼）会“松手”，更好的选择是复制上一序列的最后一个控制 $u_{T-1}$。

\begin{pitfall}[编程陷阱：仿真里代价小没下溢，一上真实尺度就 NaN]
很多人第一次自己写 MPPI，仿真里用的代价小（个位数），侥幸没遇到下溢，于是直接写 \texttt{w = exp(-S/\allowbreak lam)} 而省掉 $\rho$ 减法；一换到真实尺度的代价（几百上千）就全体下溢、归一化 $0/0$ 给出 \texttt{NaN}、机器人失控。这是新手最常见也最隐蔽的坑。\textbf{对策}：任何指数加权都先减最小（代价）/最大（奖励）值再取指数——见 \cref{eq:mppi-w-rho}。注意\emph{方向}：MPPI 用代价（最小化）减 $\min$；若实现用奖励（最大化），权重是 $e^{+R_k/\lambda}$、要减 $\max_k R_k$，搞反会把本该防下溢变成催上溢。减 $\rho$ 后仍可能有远高于最优的 rollout 权重下溢成 $0$，但这\emph{无害}——它们本就该接近零权重；$\rho$ 减法保证的是“不会所有权重一起归零”，而非“每个权重都非零”。
\end{pitfall}

\begin{practice}
\begin{enumerate}
  \item 手写一份完整 MPPI 伪代码（含 warm-start shift、$\rho$ 减法、控制约束钳位、终端代价），并对每个细节说清“不做会怎样”。
  \item 用 NumPy 把 \cref{sec:mppi-pointmass} 的一维点质量做成接收水平 MPPI：每周期采 $K$ 条带噪 rollout、算 $S_k$、按 \cref{eq:mppi-w-rho} 加权更新 $U$、执行 $u_0$、shift。观察粒子是否被拉回原点，并对比单次 \cref{eq:mppi-pm-ustar} 的闭式增益。
  \item 把时域 $T$ 从小调到大，固定 $K$，观察“远见不足”与“高维采样撞不中好序列”这两个相反失效如何在 $T$ 的两端出现。
\end{enumerate}
\end{practice}

% ════════════════════════════════════════════════════════════════
\section{采样式 MPC 家族的统一视角}
\label{sec:mppi-family}

MPPI 不是孤立的算法。把它放回采样式 MPC 的全景，会发现 CEM、iCEM、Tsallis VI-MPC、Predictive Sampling、CoVO-MPC 共享同一个数学骨架，差异都可归结为\emph{三个设计选择——用什么权重函数、用什么噪声、用什么协方差}。看清这个框架，会让你对“该怎么设计一个采样式控制器”有全局理解。

\begin{definition}[采样式 MPC 的统一骨架]\label{def:mppi-skeleton}
用一个高斯提议分布 $q=\mathcal{N}(\mu,\boldsymbol{\Sigma})$ 去逼近目标分布 $\mathbb{P}^\star\propto\mathbb{Q}\,e^{-S/\lambda}$。每一轮迭代做三步：
\begin{enumerate}
  \item \textbf{采样}：从当前 $q$ 采 $K$ 个样本 $U_1,\dots,U_K$，算各自代价 $J_1,\dots,J_K$（即前文 $S_k$）；
  \item \textbf{赋权}：给每个样本权重 $w_i=g(J_i)$，$g(\cdot)$ 是一个\emph{单调不增}的权重函数（代价越低权重越大）；
  \item \textbf{更新}：用加权样本更新分布 $\mu\leftarrow\sum_i w_i U_i/\sum_i w_i$（协方差同理）。
\end{enumerate}
\end{definition}

\noindent 这个“采样—赋权—更新”结构与机器学习的 EM 算法相通：赋权步像 E 步（按当前分布与代价给每个样本一个“重要性”），更新步像 M 步（用加权数据重估分布参数）。差别只在第二步的权重函数 $g$。为什么用\emph{高斯}作提议分布？采样便宜（计算机最擅长）、参数少且直观（$\mu,\boldsymbol{\Sigma}$ 两个量）、给定一二阶矩时熵最大（最无偏）、矩匹配有闭式。但高斯单峰、对称、形状受限——当真实最优控制分布多峰（多条等优路径）时力不从心，这是\emph{本框架所有方法的共同天花板}（突破之道如扩散模型属后续专题）。

\subsection{CEM 与 MPPI：硬阈值 vs 软加权}
\label{sec:mppi-cem}

\begin{theorem}[CEM 与 MPPI 是同一骨架的两种权重函数]\label{thm:mppi-cem}
CEM（交叉熵方法）与 MPPI 在 \cref{def:mppi-skeleton} 的第 1、3 步完全相同，只在第 2 步的权重函数不同：
\begin{equation}
  \text{MPPI:}\quad g_{\mathrm{MPPI}}(J)=e^{-J/\lambda},\qquad
  \text{CEM:}\quad g_{\mathrm{CEM}}(J)=\mathbb{1}\{J\le J_{\mathrm{elite}}\}.
  \label{eq:mppi-cem-g}
\end{equation}
MPPI 是一条\emph{光滑衰减的指数曲线}（人人有份、份额不同），CEM 是一个\emph{硬阶跃}（Heaviside 函数，代价低于精英门槛 $J_{\mathrm{elite}}$ 的样本权重 $1$、其余 $0$）。
\end{theorem}

\noindent CEM\cite{rubinstein1999cross} 起源于稀有事件估计——通过最小化提议分布与理想分布之间的\emph{交叉熵}（与 KL 密切相关，$\mathrm{KL}=$ 交叉熵 $-$ 熵）迭代改进采样分布，后被引入轨迹优化。两者“同根”在历史层面也成立：MPPI 从 KL 散度推导、CEM 从交叉熵起源，而交叉熵和 KL 本就是信息论里一对孪生量。这个统一揭示一个漂亮对应：\textbf{MPPI 的温度 $\lambda$ 与 CEM 的精英比例控制的是同一件事——筛选有多“贪婪”}。$\lambda$ 越小、指数曲线越陡、权重越集中到极少数最低代价样本（越贪婪、越像 CEM 取很小精英比例，$\lambda\to0$ 退化为只取最好的一个）；$\lambda$ 越大、曲线越平、所有样本趋于等权（越保守）。CEM 精英比例越小越像小 $\lambda$ 的 MPPI、越大越像大 $\lambda$ 的 MPPI——它们是\emph{同一个“探索-利用”旋钮的两种刻度}。各自的“脾气”：CEM 直观、无需调温度（精英比例不依赖代价数值尺度），但信息利用粗糙（精英内不分高下、门槛外哪怕只差一点也彻底扔掉），且精英比例取太小则方差估计噪声大、易过早坍缩；MPPI 信息利用充分（所有样本按代价平滑赋权），但 $\lambda$ 难调且依赖代价尺度（代价尺度一变就要重调 $\lambda$）。这对“硬阈值 vs 软加权”在深度学习里有一个孪生兄弟——hard attention vs soft attention（top-k vs softmax），取舍完全平行。

\subsection{iCEM、Tsallis、Predictive Sampling、CoVO-MPC}
\label{sec:mppi-variants}

\paragraph{iCEM：让 CEM 快到能实时。}原始 CEM 的硬伤是\emph{慢}：它要每个控制周期内迭代多轮，而每轮又要大量样本才能让高斯估计稳定，且\emph{轮间串行}（这一轮的分布由上一轮的精英决定）——GPU 能压缩“每轮时间”却压不了“轮数”，于是无法实时。iCEM（Pinneri 等，CoRL 2020）用四个改进把样本数砍到能实时：原文摘要报告\emph{所需样本减少 $2.7$--$22\times$、性能提升 $1.2$--$10\times$}\cite{pinneri2020sample}（“$13.7\times$ 更少样本达 $90\%$ 成功率”等更细的单任务数字为源 md 转述、未在摘要核到，此处略）：(1) \emph{有色噪声}——不再对每时刻独立采白噪声，而是采功率谱按 $S(f)\propto1/f^\beta$ 衰减的有色噪声（$\beta=2$ 布朗噪声对控制序列最友好，生成的扰动平滑、时间相关，更可能成精英），是影响最大的一招；(2) \emph{elite memory}——保留上一轮精英加进下一轮采样池，但旧精英要重新评估、过时的自然淘汰（与 RL 的 replay buffer 相通）；(3) \emph{shift init}——把上一控制周期的解整体左移一步作为新周期初始均值（warm-start 在 CEM 上的应用）；(4) \emph{mean-as-sample}——把当前分布的均值本身作为额外样本加入评估，补救高维高斯的“薄壳现象”（$d$ 维高斯典型样本离均值约 $\sigma\sqrt d$，球心附近几乎采不到，而均值常是很优的点）。这四招在 MPPI 世界都有对应（有色噪声 $\leftrightarrow$ Smooth-MPPI、shift init $\leftrightarrow$ warm-start），再次印证“同根”。

\begin{theorem}[Tsallis VI-MPC：一个参数连续统一 MPPI 与 CEM]\label{thm:mppi-tsallis}
用\emph{变形指数}（deformed / q-exponential）$\exp_r(x)=[\,1+(1-r)\,x\,]_+^{\frac{1}{1-r}}$（$[\cdot]_+=\max(\cdot,0)$）作权重函数 $w_i=\exp_r(-J_i/\lambda)$，则参数 $r$ 是一个连续的“筛选硬度”旋钮：$r\to1$ 时 $\exp_r\to\exp$、退化为 MPPI 的软指数加权；$r\to\infty$ 时趋向 CEM 的硬阈值（精英内近等权、非精英为零）。MPPI 与 CEM 是这条由 $r$ 参数化的连续谱的\emph{两个端点}\cite{wang2021variational}。
\end{theorem}

\noindent 变形指数来自非广延统计力学（Tsallis 1988\cite{tsallis1988possible}）：标准玻尔兹曼分布的指数 $e^{-E/kT}$ 在系统“非广延”（熵不可加）时被推广为 $\exp_r$，$r$（物理里常记 $q$）刻画偏离广延性的程度、$r=1$ 恢复标准情形。这解释了 $r=1$ 为何特殊。更深一层：权重函数的形状根源于\emph{用什么散度衡量分布的接近程度}——MPPI 用 KL 散度得到标准指数权重，Tsallis 把 KL 换成更一般的\emph{Tsallis 散度}（$r=1$ 退化为 KL），对应的最优权重就从标准指数变成变形指数。这把 \cref{thm:mppi-cem} 的本质洞察从“两种权重函数”升级为“两种散度度量”，统一表见 \cref{tab:mppi-tsallis}。

\begin{table}[H]\centering\small
\caption{Tsallis 框架统一 MPPI 与 CEM}
\label{tab:mppi-tsallis}
\begin{tabular}{llll}
\toprule
维度 & MPPI & CEM & Tsallis VI-MPC \\
\midrule
散度正则 & KL 散度 & （隐式 Heaviside 筛选） & Tsallis 散度（KL 的推广） \\
权重函数 & $e^{-J/\lambda}$（标准指数） & $\mathbb{1}\{J\le J_{\mathrm{elite}}\}$（硬阶跃） & $\exp_r(-J/\lambda)$（变形指数） \\
参数 $r$ & $r\to1$ & $r\to\infty$ & 可调 $1<r<\infty$ \\
筛选硬度 & 最软（人人有份） & 最硬（非精英即弃） & 连续可调 \\
\bottomrule
\end{tabular}
\end{table}

\paragraph{Predictive Sampling：极简端点。}Predictive Sampling（Howell 等，2022，MuJoCo MPC）把“赋权”这一步彻底砍掉——不做指数加权、不做精英筛选，就一句话：\emph{采一批样本，取代价最低的那一个作为新均值}（取 $\arg\min$），没有温度、没有精英比例、没有变形参数\cite{howell2022predictive}。它本是 MJPC 里一个“最朴素的 baseline、主要为教学价值”，结果竟与更成熟的 iLQG、梯度下降竞争力相当（原文摘要语\cite{howell2022predictive}）。为何意外有效？\emph{模型可用且 rollout 便宜}（MuJoCo 快速仿真）、\emph{样条参数化降维}（用少数样条控制点参数化整条控制序列，把高维逐时刻控制压成低维，$\arg\min$ 在低维空间撞中好样本的概率大增——呼应薄壳现象）、\emph{任务确定性}（噪声小，优胜者诅咒不严重）。它在 \cref{thm:mppi-tsallis} 框架里是\emph{最硬端点}：$\arg\min$ 等价于 MPPI 在 $\lambda\to0$、CEM 在精英数 $=1$、Tsallis 在 $r\to\infty$ 的极限。它的弱点是\emph{优胜者诅咒（winner's curse）}——代价评估带噪声时，$\arg\min$ 系统性地选中“真实代价不错、且恰好噪声为负（运气好）”的样本，高噪声下被带偏；MPPI 的加权平均天然抵抗（单样本噪声在平均中被抵消）。

\begin{theorem}[CoVO-MPC：MPPI 的收敛性与最优协方差]\label{thm:mppi-covo}
考虑代价是二次型的情形（$J(U)$ 关于控制序列 $U$ 是二次函数，涵盖时变 LQR）。CoVO-MPC（Yi、Pan 等，L4DC 2024）证明 MPPI 的迭代满足
\begin{equation}
  \|U_{k+1}-U^\star\|\le\rho(\boldsymbol{\Sigma})\,\|U_k-U^\star\|,\qquad 0\le\rho<1,
  \label{eq:mppi-covo}
\end{equation}
即\emph{线性（几何）收敛}到最优控制序列 $U^\star$。收缩率 $\rho(\boldsymbol{\Sigma})$ 由采样协方差 $\boldsymbol{\Sigma}$、温度 $\lambda$、代价的 Hessian 共同决定；求让 $\rho$ 最小的 $\boldsymbol{\Sigma}^\star$，得最优采样协方差，其形状由代价的\emph{局部 Hessian} 决定，遵循“凸方向（曲率大）小方差、平方向（曲率小）大方差”的原则。据此的 CoVO-MPC 算法比 vanilla MPPI 提升 $43$--$54\%$\cite{yi2024covo}。
\end{theorem}

\noindent 这是 MPPI 的首个收敛性分析与最优协方差设计（原文定位\cite{yi2024covo}）。它解释了一个重要实践现象：\emph{几何收敛意味着收敛很快}，所以 MPPI 通常每控制周期只迭代很少几次（甚至 $1$ 次）就够用——配合 warm-start（每周期从上次的好解出发），一次迭代就能跟上。为何不直接用牛顿法？因为牛顿法要求代价可微、能解析算梯度和 Hessian，而采样式 MPC 的存在理由正是处理\emph{不可微、黑箱}的代价和动力学；CoVO-MPC 用采样和 rollout \emph{数值估计} Hessian，把“用 Hessian 加速”的二阶好处移植进无导数的采样框架——它站在采样式方法（能处理黑箱）与二阶方法（收敛快）的交叉点上。注意 $\boldsymbol{\Sigma}$ 与 $\lambda$ 在 $\rho$ 里耦合：$\lambda$ 管“全局探索强度”、$\boldsymbol{\Sigma}$ 的形状管“探索强度在各方向如何分配”，求最优 $\boldsymbol{\Sigma}$ 是在给定 $\lambda$ 下求的——所以改了 $\boldsymbol{\Sigma}$ 可能需要重新审视 $\lambda$。

\begin{insight}[一条完整的“筛选硬度”谱：形状比数目更重要]
把本节方法挂到同一张框架图上，会看到一条完整的谱：从“几乎不筛选”（$\lambda$ 很大、权重近等权）到“筛选到极致”（$\arg\min$、只取一个）。MPPI（$r\to1$）、Tsallis（中间 $r$）、CEM（$r\to\infty$）、Predictive Sampling（最硬端点）全是这条谱上的点，差异只在\emph{权重函数（$\lambda$/精英比例/$r$）}这一个维度；iCEM 改的是\emph{噪声}（有色噪声）这个维度，CoVO-MPC 改的是\emph{协方差}（贴合 Hessian）这个维度。三个维度合起来就是 \cref{def:mppi-skeleton} 的全部设计空间。这条谱最深刻的启示有二：其一，\emph{采样式 MPC 的有效性主要来自“采样 + 用模型评估真实代价”这个核心机制，权重方案只是在这个核心之上调节探索-利用的平衡}（这解释了为何连最极端的 $\arg\min$ 都能 work）；其二，CoVO-MPC 点破了“\emph{采样分布的形状（协方差）比采样数目更重要}”——与其盲目加大 $K$，不如把协方差贴合代价地形的曲率。面对一个新任务，先用最简单的（甚至 Predictive Sampling）跑通探路，再按需精化权重/噪声/协方差。
\end{insight}

\begin{pitfall}[思维陷阱：把“硬阈值 vs 软加权”看成非此即彼的二元对立]
认为一个采样式 MPC 要么是“CEM 派”（硬阈值）、要么是“MPPI 派”（软加权），二者泾渭分明。\textbf{后果}：遇到一个“既不完全硬阈值、也不完全指数”的方法（Tsallis VI-MPC、某些 top-k 加权变体）时无法归类，觉得它“四不像”。\textbf{根本原因}：硬阈值和指数加权其实是一条\emph{连续谱}的两端，而非两个离散点——Heaviside 阶跃是某极限下的特例（\cref{thm:mppi-tsallis}），中间有无数“软硬程度不同”的权重函数。\textbf{对策}：把权重函数的“硬度”看成一个连续可调的量（$\exp_r$ 的 $r$）；遇到任何采样式 MPC，先问“它的权重函数长什么样、有多硬、改的是三维框架里哪个维度（权重/噪声/协方差）”，而非急着塞进 CEM 或 MPPI 的盒子。
\end{pitfall}

\begin{practice}
\begin{enumerate}
  \item 基于 \cref{def:mppi-skeleton} 的三步骨架，用 NumPy 实现一个\emph{可切换权重函数}的采样式优化器：传入 \texttt{weight\_fn(J)}，分别用指数权重 $e^{-J/\lambda}$、硬阈值 $\mathbb{1}\{J\le J_{\mathrm{elite}}\}$、变形指数 $\exp_r$、$\arg\min$ 四种，在同一到达任务上对比收敛曲线。验证 MPPI 的 $\lambda$ 调小、CEM 的精英比例调小是否都让收敛更快但在多峰任务上更易陷次优。
  \item （统一框架定位）论证 $\arg\min$ 等价于 MPPI 在 $\lambda\to0$、CEM 在精英数 $=1$、Tsallis 在 $r\to\infty$ 的极限，把 Predictive Sampling 准确标到 \cref{tab:mppi-tsallis}。
  \item 实现一个有色噪声生成器（白噪声频谱 $\times\,1/f^{\beta/2}$ 整形 $\to$ 逆 FFT $\to$ 按行归一化到单位标准差），扫描 $\beta\in\{0,1,2,3\}$，验证 $\beta=2$ 附近是甜点、且“样本越少有色噪声优势越大”。再把它迁移到 MPPI 的采样上，对比控制平滑度。
\end{enumerate}
\end{practice}

% ════════════════════════════════════════════════════════════════
\section*{本章常见误解汇总}
\begin{table}[H]\centering\small
\begin{tabular}{p{0.46\textwidth} p{0.46\textwidth}}
\toprule
\textbf{常见误解} & \textbf{正确理解} \\
\midrule
指数变换对任意最优控制问题都能线性化 & 仅在 Kappen 条件 $\boldsymbol{\Sigma}=\lambda R^{-1}$（控制仿射 + 二次控制代价 + 噪声进控制通道）下相消（\cref{thm:mppi-kappen}）\\
$\boldsymbol{\Sigma},R,\lambda$ 可独立设定、事后微调 & $\boldsymbol{\Sigma}=\lambda R^{-1}$ 是充要条件，两自由一被定（\cref{eq:mppi-kappen}）\\
desirability $\psi$ 是概率密度 & $\psi=e^{-J/\lambda}$ 是非归一化的合意度场，相对大小有意义（\cref{thm:mppi-kappen}）\\
Feynman--Kac 直接给出最优轨迹 & 给的是标量场 $\psi$，最优控制是 $\nabla\log\psi$ 这个反馈律（\cref{thm:mppi-control-law}）\\
权重里的代价 $S_k$ 必须可微 & $S_k$ 只作标量进指数、从不求导；可不连续/黑箱（\cref{sec:mppi-nondiff}）\\
MPPI 更新是“梯度下降走一步” & 是闭式吉布斯分布的加权平均；无学习率（\cref{thm:mppi-weights} 洞见）\\
$K$ 越大一定越好 & $K$ 只降方差；提议分布离 $\mathbb{P}^\star$ 太远则白采（\cref{sec:mppi-mc}）\\
温度 $\lambda$ 与协方差 $\boldsymbol{\Sigma}$ 作用相同 & $\boldsymbol{\Sigma}$ 管“撒哪里”、$\lambda$ 管“信哪条”（\cref{sec:mppi-temperature}）\\
随机环境里 $\lambda$ 越小越精准 & $\lambda$ 同时是风险旋钮，越小越冒险、越脆弱（\cref{thm:mppi-meanvar}）\\
不减 $\rho$ 直接 $e^{-S/\lambda}$ 也能跑 & 真实代价尺度下全体下溢成 NaN（\cref{thm:mppi-rho}）\\
MPPI 与 CEM 是两个孤立算法 & 同一“采样—赋权—更新”骨架的两种权重函数（\cref{thm:mppi-cem}）\\
硬阈值 vs 软加权非此即彼 & 是 $\exp_r$ 的 $r$ 连续参数化的一条谱的两端（\cref{thm:mppi-tsallis}）\\
加大 $K$ 是提升性能的主路 & 协方差形状比采样数目更重要（\cref{thm:mppi-covo} 洞见）\\
MPPI 是孤立控制技巧，与 RL 无关 & 与 REINFORCE/SAC/MPO 同为控制即推断的实例（\cref{sec:mppi-inference}）\\
\bottomrule
\end{tabular}
\end{table}

% ════════════════════════════════════════════════════════════════
\section*{本章小结}

\begin{table}[H]\centering\small
\caption{符号表}
\begin{tabular}{lll}
\toprule
符号 & 含义 & 首次出现 \\
\midrule
$f(x),G(x)$ & 漂移 / 控制矩阵（控制与噪声同通道 $G$） & \cref{eq:mppi-sde} \\
$\boldsymbol{\Sigma},R,\lambda$ & 噪声/采样协方差 / 控制代价权重 / 温度 & \cref{def:mppi-problem} \\
$J(x,t),\ \psi=e^{-J/\lambda}$ & 值函数 / desirability（合意度，非概率） & \cref{eq:mppi-cost,thm:mppi-kappen} \\
$\boldsymbol{\Sigma}=\lambda R^{-1}$ & Kappen 条件（线性化充要） & \cref{eq:mppi-kappen} \\
$V(x,t),\phi(x_T)$ & 运行代价 / 终端代价 & \cref{eq:mppi-cost} \\
$\mathbb{Q},\mathbb{P}^\star$ & 基/提议分布 / 吉布斯最优轨迹分布 & \cref{def:mppi-freeenergy,eq:mppi-gibbs} \\
$S(\tau)=S_k$ & 一条轨迹/rollout 的总代价（标量、可黑箱） & \cref{eq:mppi-S} \\
$\mathcal{F}(\mathbb{Q})=-\lambda\log Z$ & 自由能 / 配分函数 $Z=\mathbb{E}_{\mathbb{Q}}[e^{-S/\lambda}]$ & \cref{eq:mppi-F} \\
$w_k,\rho=\min_k S_k$ & 归一化权重 / 数值稳定基准 & \cref{eq:mppi-w,eq:mppi-w-rho} \\
$\boldsymbol{\varepsilon}_k,U,T$ & 第 $k$ 条噪声序列 / 控制序列 / 时域 & \cref{thm:mppi-weights} \\
$\mathrm{ESS}=1/\sum_k w_k^2$ & 有效样本数（$\in[1,K]$，调 $\lambda$ 的探针） & \cref{eq:mppi-ess} \\
$z=e^{-v/\lambda},\ z=MPz$ & LMDP desirability / 线性贝尔曼方程 & \cref{eq:mppi-zmpz} \\
$\mathcal{O}_t$ & 最优性变量（控制即推断） & \cref{def:mppi-inference} \\
$\exp_r,r$ & Tsallis 变形指数 / 筛选硬度参数 & \cref{thm:mppi-tsallis} \\
$\rho(\boldsymbol{\Sigma})$ & CoVO-MPC 收缩率（局部，$\ne\min_k S_k$） & \cref{eq:mppi-covo} \\
\bottomrule
\end{tabular}
\end{table}

\begin{table}[H]\centering\small
\caption{定理速查表}
\begin{tabular}{lll}
\toprule
定理 / 公式 & 一句话说明 & 对应 \\
\midrule
Kappen 指数变换 \cref{eq:mppi-linearpde} & $\psi=e^{-J/\lambda}$ 把非线性 HJB 整体线性化 & \cref{thm:mppi-kappen} \\
最优控制 \cref{eq:mppi-ustar} & $u^\star=-R^{-1}G^\top\nabla J$，正比值函数负梯度 & \cref{thm:mppi-optu} \\
Feynman--Kac \cref{eq:mppi-fk} & 线性 PDE 解 = 无控轨迹上的期望 & \cref{thm:mppi-feynmankac} \\
路径积分控制律 \cref{eq:mppi-pi-control} & $u^\star=\lambda R^{-1}G^\top\nabla\log\psi$，采样平均 & \cref{thm:mppi-control-law} \\
路径积分 = LQR \cref{thm:mppi-lqr} & LQ 特例逐字等于 Riccati 解 & \cref{sec:mppi-pointmass} \\
线性贝尔曼 \cref{eq:mppi-zmpz} & 离散镜像 $z=MPz$，KL 控制代价 & \cref{thm:mppi-lmdp} \\
自由能--KL 对偶 \cref{eq:mppi-dual} & $\mathcal{F}=\min_{\mathbb{P}}\{\mathbb{E}[S]+\lambda\mathrm{KL}\}$，最优吉布斯 & \cref{thm:mppi-duality} \\
MPPI 权重 \cref{eq:mppi-w} & $w_k\propto e^{-S_k/\lambda}$，自归一化重要性权重 & \cref{thm:mppi-weights} \\
$\rho$ 减法 \cref{eq:mppi-w-rho} & log-sum-exp，必需的数值保命 & \cref{thm:mppi-rho} \\
MPPI $\approx$ REINFORCE & 同骨架，指数权重 vs 线性权重 & \cref{tab:mppi-reinforce} \\
均值-方差展开 \cref{eq:mppi-meanvar} & $\mathcal{F}\approx\mathbb{E}[S]-\mathrm{Var}[S]/2\lambda$，$\lambda$ 是风险旋钮 & \cref{thm:mppi-meanvar} \\
规划即推断 \cref{eq:mppi-posterior} & $\mathbb{P}^\star$ = 全程最优证据下的后验 & \cref{def:mppi-inference} \\
CEM = MPPI 换权重 \cref{eq:mppi-cem-g} & 软指数 vs 硬阈值，$\lambda\leftrightarrow$ 精英比例 & \cref{thm:mppi-cem} \\
Tsallis 统一 & $\exp_r$ 的 $r$ 连续插值 MPPI($r{\to}1$)/CEM($r{\to}\infty$) & \cref{thm:mppi-tsallis} \\
CoVO-MPC 收敛 \cref{eq:mppi-covo} & 二次代价下线性收敛；最优 $\boldsymbol{\Sigma}$ 由 Hessian 定 & \cref{thm:mppi-covo} \\
\bottomrule
\end{tabular}
\end{table}

\begin{table}[H]\centering\small
\caption{知识点总表}
\begin{tabular}{clll}
\toprule
\# & 知识点 & 核心要点 & 难度 \\
\midrule
1 & 指数变换线性化 HJB & $\psi=e^{-J/\lambda}$；Kappen 条件相消非线性项 & \(\bigstar\bigstar\bigstar\) \\
2 & Feynman--Kac 与控制律 & 线性 PDE 解成期望；最优控制采样平均出来 & \(\bigstar\bigstar\bigstar\) \\
3 & 路径积分 = LQR & 一维点质量蒙特卡洛逐字等于 Riccati & \(\bigstar\bigstar\bigstar\) \\
4 & 线性可解 MDP & 离散镜像 $z=MPz$；KL 控制代价；可叠加 & \(\bigstar\bigstar\bigstar\bigstar\) \\
5 & 自由能--KL 对偶 & Gibbs 变分原理；脱离 Kappen 假设 & \(\bigstar\bigstar\bigstar\bigstar\) \\
6 & MPPI 权重公式 & $w_k\propto e^{-S_k/\lambda}$；$\rho$ 减法；高斯密度消去 & \(\bigstar\bigstar\bigstar\) \\
7 & MPPI $\approx$ REINFORCE & 零阶策略梯度闭式特例；自然梯度/无学习率 & \(\bigstar\bigstar\bigstar\bigstar\) \\
8 & 控制即推断 & MPPI/SAC/MPO/AWR 同一后验，按资源选路线 & \(\bigstar\bigstar\bigstar\bigstar\) \\
9 & 温度 $\lambda$ 三张面孔 & 噪声温度 / softmax 锐度 / 风险态度 & \(\bigstar\bigstar\bigstar\) \\
10 & 风险敏感 & 自由能 = 指数效用；低温在随机环境冒险 & \(\bigstar\bigstar\bigstar\bigstar\) \\
11 & 接收水平 MPPI & 开环序列 + 重规划 = 闭环；warm-start 即好提议 & \(\bigstar\bigstar\) \\
12 & 采样式 MPC 统一骨架 & 采样—赋权—更新；权重/噪声/协方差三维 & \(\bigstar\bigstar\bigstar\) \\
13 & CEM 与 Tsallis & 软硬两端；$\exp_r$ 连续谱 & \(\bigstar\bigstar\bigstar\bigstar\) \\
14 & iCEM 与 Predictive Sampling & 样本高效四招；极简 $\arg\min$ 端点 & \(\bigstar\bigstar\bigstar\) \\
15 & CoVO-MPC 收敛 & 线性收敛；形状比数目更重要 & \(\bigstar\bigstar\bigstar\bigstar\) \\
\bottomrule
\end{tabular}
\end{table}

% ════════════════════════════════════════════════════════════════
\section*{延伸阅读}
\begin{itemize}
  \item \textbf{理论起源}：Kappen《Linear theory for control of nonlinear stochastic systems》\cite{kappen2005linear} 与《Path integrals and symmetry breaking for optimal control theory》\cite{kappen2005path}——指数变换、Kappen 条件、双缝对称性破缺的奠基作；离散镜像见 Todorov《Efficient computation of optimal actions》\cite{todorov2009efficient} 与可组合性\cite{todorov2009compositionality}。
  \item \textbf{信息论 MPPI（主线）}：Williams 等《Model Predictive Path Integral Control: From Theory to Parallel Computation》\cite{williams2017model}（JGCD，generalized importance sampling + GPU，MPPI 诞生）与《Information-Theoretic Model Predictive Control: Theory and Applications to Autonomous Driving》\cite{williams2018information}（T-RO，自由能--KL 完整推导，本章 \cref{sec:mppi-freeenergy} 的主要出处）；策略搜索前身 PI$^2$ 见 Theodorou 等\cite{theodorou2010generalized}。
  \item \textbf{控制即推断与风险敏感}：Levine《Reinforcement learning and control as probabilistic inference》\cite{levine2018reinforcement}（统一视角综述）、Kappen--Gómez--Opper\cite{kappen2012optimal}（KL 控制含路径积分为特例）；风险敏感见 van den Broek 等\cite{vandenbroek2010risk} 与 Jacobson 的 LEQG\cite{jacobson1973optimal}。
  \item \textbf{家族与统一框架}：iCEM\cite{pinneri2020sample}、Tsallis VI-MPC\cite{wang2021variational}、Predictive Sampling / MuJoCo MPC\cite{howell2022predictive}、CoVO-MPC\cite{yi2024covo}；最大熵 RL 一侧见 SAC\cite{haarnoja2018soft}、MPO\cite{abdolmaleki2018maximum}、AWR\cite{peng2019advantage}。
  \item \textbf{数学工具}：Feynman--Kac 的母体 Kac 1949\cite{kac1949distributions} 与 Feynman 路径积分\cite{feynman1948spacetime}；Cole--Hopf 变换 Hopf 1950\cite{hopf1950partial}、Cole 1951\cite{cole1951quasi}；策略梯度 REINFORCE\cite{williams1992simple}。
\end{itemize}

\section*{本章与后续章节的关系}
\begin{table}[H]\centering\small
\begin{tabular}{lll}
\toprule
后续章节 & 关系 & 本章铺垫 \\
\midrule
\cref{sec:plan-traj} 轨迹优化 & 黑箱代价的随机版替代 & 代价泛函 \cref{eq:mppi-S}、不可微代价 \\
梯度式 MPC（iLQR/DDP） & 互补：梯度法可微局部 vs MPPI 黑箱全局 & \cref{thm:mppi-covo} 的二阶/采样交叉点 \\
最大熵 RL（SAC/MPO） & 同一 KL 正则推断的离线摊销版 & 控制即推断（\cref{sec:mppi-inference}）\\
基于模型的 RL（TD-MPC 等） & 采样式规划嵌入世界模型 & 统一骨架（\cref{def:mppi-skeleton}）、iCEM 样本效率 \\
不确定性 / 安全规划 & 风险敏感是最平滑入口 & 自由能 = 指数效用（\cref{thm:mppi-meanvar}）\\
扩散 / 生成式规划 & 突破高斯提议分布的表达天花板 & 提议分布质量主线（\cref{sec:mppi-family}）\\
\bottomrule
\end{tabular}
\end{table}

\section*{故障排查手册}
\begin{enumerate}
  \item \textbf{症状}：控制输出 \texttt{NaN}、机器人失控。\textbf{原因}：未做 $\rho$ 减法，真实代价尺度下 $e^{-S/\lambda}$ 全体下溢、归一化 $0/0$（\cref{thm:mppi-rho}）。\textbf{对策}：权重一律按 \cref{eq:mppi-w-rho} 减 $\min_k S_k$（奖励表述减 $\max$）。
  \item \textbf{症状}：控制信号高频抖动、电机嗡鸣。\textbf{原因}：原始 MPPI 输出是随机轨迹加权平均、天然带高频噪声。\textbf{对策}：加 Savitzky--Golay 平滑（\cref{alg:mppi} 第 4 步）；或用有色噪声采样（\cref{sec:mppi-variants}）。
  \item \textbf{症状}：加大 $K$ 也不收敛、$\mathrm{ESS}$ 长期接近 $1$。\textbf{原因}：提议分布离 $\mathbb{P}^\star$ 太远，低代价区无样本落入（\cref{sec:mppi-mc}）。\textbf{对策}：改善 warm-start、调大/贴合 $\boldsymbol{\Sigma}$、注入先验；分布质量优先于样本数。
  \item \textbf{症状}：权重坍缩到一两条 rollout、控制剧烈抖动。\textbf{原因}：$\lambda$ 太小、$\mathrm{ESS}\to1$（\cref{thm:mppi-lambda}）。\textbf{对策}：用 ESS 当探针调大 $\lambda$，或按目标 ESS 自动调温。
  \item \textbf{症状}：随机/不确定环境里控制激进而脆弱。\textbf{原因}：$\lambda$ 太小，落入 risk-seeking（\cref{thm:mppi-meanvar}）。\textbf{对策}：理解 $\lambda$ 为风险态度、需保守时不一味降温；安全攸关任务转向 risk-averse 推广。
  \item \textbf{症状}：换个机器人/改了代价权重后 $\lambda$ 失效。\textbf{原因}：$\lambda$ 依赖代价数值尺度（\cref{thm:mppi-cem}）。\textbf{对策}：用 ESS 重新标定 $\lambda$，或改用对尺度不敏感的 CEM 精英比例 / 把控制代价吸收进 $\lambda$。
\end{enumerate}
